\documentclass[11pt]{article}
\usepackage[margin=1in]{geometry}
\usepackage[T1]{fontenc}
\usepackage[utf8]{inputenc}
\usepackage[french]{babel}
\usepackage{amsmath,amssymb,amsthm}
\usepackage[colorlinks=true,linkcolor=blue,urlcolor=blue]{hyperref}
\newtheorem{problem}{Problème}
\newenvironment{refs}{\par\small\noindent\emph{Références :}\begin{itemize}\setlength\itemsep{0pt}}{\end{itemize}\normalsize}
\title{Hors de la Caverne \\ \Large Analyse fonctionnelle}
\author{}
\date{}
\begin{document}
\maketitle
\noindent\emph{Des problèmes, pas de réponses.} Chaque problème ci-dessous est posé
sans solution. Les références indiquent où trouver les connaissances nécessaires.
\medskip


\section*{Licence}

\begin{problem}[{Deux normes sur C[0,1] --- Licence}]
\textbf{FA-01.} \textbf{Problème.} Soit $ C[0,1] $ l'espace vectoriel des fonctions continues à valeurs réelles
sur $ [0,1] $. Munissez-le de la norme sup puis, séparément, de la norme intégrale :
\[ \|f\|_\infty = \sup\{\,|f(x)| : 0 \le x \le 1\,\}, \qquad \|f\|_1 = \int_0^1 |f(x)|\,dx . \]
\begin{enumerate}
\item[(a)] Démontrez que $ (C[0,1], \|\cdot\|_\infty) $ est complet : toute suite de Cauchy converge vers
une fonction continue.
\item[(b)] Exhibez une suite de $ C[0,1] $ qui est de Cauchy pour $ \|\cdot\|_1 $ mais n'a pas de limite
dans $ C[0,1] $, et concluez que $ (C[0,1], \|\cdot\|_1) $ n'est pas complet.
\item[(c)] Montrez que les deux normes ne sont pas équivalentes : il n'existe aucune constante $ K $ telle
que $ \|f\|_\infty \le K\|f\|_1 $ pour toute $ f $.
\end{enumerate}

\end{problem}

\noindent En dimension finie toutes les normes sont équivalentes et tout espace normé est
complet, si bien que rien de tout cela ne peut s'y voir. Le défaut de complétude de la question (b)
n'est pas un défaut de la norme intégrale mais de l'espace : compléter $ C[0,1] $ pour
$ \|\cdot\|_1 $ est précisément ce qui produit l'espace de Lebesgue $ L^1[0,1] $, découverte
(Riesz et Fischer, 1907) qui a rendu l'intégrale de Lebesgue indispensable à l'analyse. Ce problème
est la porte d'entrée de toute la discipline : la norme que vous choisissez détermine quelles limites
existent.

\begin{refs}
\item Contexte : Espace complet --- Wikipédia (\url{https://fr.wikipedia.org/wiki/Espace_complet})
\item Contexte : Convergence uniforme --- Wikipédia (\url{https://fr.wikipedia.org/wiki/Convergence_uniforme}), Espace Lp --- Wikipédia (\url{https://fr.wikipedia.org/wiki/Espace_Lp})
\item Lecture : E. Kreyszig, \textit{Introductory Functional Analysis with Applications}, ch. 1--2
\item Cours : MIT OCW 18.102 Introduction to Functional Analysis (en anglais) (\url{https://ocw.mit.edu/courses/18-102-introduction-to-functional-analysis-spring-2021/})
\end{refs}

\bigskip

\begin{problem}[{Ce qui fait un espace de Banach --- Licence}]
\textbf{FA-02.} \textbf{Problème.} Un \textit{espace de Banach} est un espace vectoriel normé complet pour la distance
induite par sa norme.
\begin{enumerate}
\item[(a)] Démontrez qu'un espace normé $ X $ est un espace de Banach si et seulement si toute série
absolument convergente de $ X $ converge : dès que $ \sum_n \|x_n\| < \infty $, les sommes
partielles de $ \sum_n x_n $ convergent dans $ X $.
\item[(b)] À l'aide de ce critère ou directement, démontrez que les espaces de suites $ \ell^p $ pour
$ 1 \le p < \infty $, munis de la norme $ \|x\|_p = \big(\sum_n |x_n|^p\big)^{1/p} $, et
$ \ell^\infty $ (suites bornées avec la norme sup) sont des espaces de Banach.
\end{enumerate}

\end{problem}

\noindent Stefan Banach a axiomatisé ces espaces dans sa thèse de doctorat de 1920 à Lwów, et
la théorie a grandi autour du Café écossais, où Banach, Steinhaus, Mazur et Ulam échangeaient des
problèmes dans un cahier qui nous est parvenu. Le critère de séries de la question (a) est le cheval
de trait déguisé en curiosité : c'est ainsi que l'on démontre la complétude des espaces quotients, de
$ L^p $ et des espaces d'opérateurs de FA-06, et c'est l'ombre exacte, en dimension infinie, du fait
que la convergence absolue entraîne la convergence dans $\mathbb{R}$.

\begin{refs}
\item Contexte : Espace de Banach --- Wikipédia (\url{https://fr.wikipedia.org/wiki/Espace_de_Banach})
\item Contexte : Espace de suites $\ell$p --- Wikipédia (\url{https://fr.wikipedia.org/wiki/Espace_de_suites_%E2%84%93p})
\item Lecture : E. Kreyszig, \textit{Introductory Functional Analysis with Applications}, ch. 2
\item Lecture : J. B. Conway, \textit{A Course in Functional Analysis}, ch. III
\end{refs}

\bigskip

\begin{problem}[{Géométrie de l'espace de Hilbert : le point le plus proche --- Licence}]
\textbf{FA-03.} \textbf{Problème.} Soit $ H $ un espace de Hilbert --- un espace vectoriel muni d'un produit
scalaire $ \langle\cdot,\cdot\rangle $ et complet pour la norme
$ \|x\| = \sqrt{\langle x, x\rangle} $.
\begin{enumerate}
\item[(a)] Démontrez l'identité du parallélogramme
\[ \|x+y\|^2 + \|x-y\|^2 = 2\|x\|^2 + 2\|y\|^2 . \]
\item[(b)] Démontrez que tout convexe fermé non vide $ C \subset H $ contient un unique point de norme
minimale, et plus généralement un unique point le plus proche d'un $ x \in H $ donné.
\item[(c)] Déduisez-en le théorème de projection : pour tout sous-espace fermé $ M $, tout $ x \in H $
se décompose de manière unique en $ x = m + n $ avec $ m \in M $ et $ n \perp M $, de sorte
que $ H = M \oplus M^{\perp} $.
\item[(d)] Utilisez (c) pour démontrer le théorème de représentation de Riesz : toute forme linéaire bornée
$ \varphi $ sur $ H $ s'écrit $ \varphi(x) = \langle x, y\rangle $ pour un unique
$ y \in H $, avec $ \|\varphi\| = \|y\| $.
\end{enumerate}

\end{problem}

\noindent Cet enchaînement --- identité du parallélogramme, point le plus proche, décomposition
orthogonale, dualité --- est la raison pour laquelle l'espace de Hilbert est le seul espace de dimension
infinie où l'intuition euclidienne survit à peu près intacte. David Hilbert a introduit l'espace
$ \ell^2 $ vers 1906 en étudiant les équations intégrales ; von Neumann en a donné les axiomes
abstraits en 1929, et Frigyes Riesz (ainsi que Fréchet, indépendamment) avait déjà trouvé le théorème
de représentation en 1907. Toute la mécanique quantique et toute la théorie $ L^2 $ des équations
aux dérivées partielles reposent sur la question (c).

\begin{refs}
\item Contexte : Espace de Hilbert --- Wikipédia (\url{https://fr.wikipedia.org/wiki/Espace_de_Hilbert})
\item Contexte : Théorème de représentation de Riesz --- Wikipédia (\url{https://fr.wikipedia.org/wiki/Th%C3%A9or%C3%A8me_de_repr%C3%A9sentation_de_Riesz_(Fr%C3%A9chet-Riesz)})
\item Lecture : E. Kreyszig, \textit{Introductory Functional Analysis with Applications}, ch. 3
\item Lecture : H. Brezis, \textit{Analyse fonctionnelle : théorie et applications}, ch. 5
\end{refs}

\bigskip

\begin{problem}[{Bases orthonormées et l'unique espace de Hilbert --- Licence}]
\textbf{FA-04.} \textbf{Problème.} Soit $ H $ un espace de Hilbert et $ (e_n) $ une suite orthonormée de
$ H $.
\begin{enumerate}
\item[(a)] Démontrez l'inégalité de Bessel : $ \sum_n |\langle x, e_n\rangle|^2 \le \|x\|^2 $ pour tout
$ x \in H $.
\item[(b)] On dit qu'une famille orthonormée est \textit{totale} (une base orthonormée) si son sous-espace
engendré fermé est $ H $. Démontrez que $ (e_n) $ est totale si et seulement si l'identité de
Parseval
\[ \sum_n |\langle x, e_n\rangle|^2 = \|x\|^2 \]
vaut pour tout $ x $.
\item[(c)] Démontrez que tout espace de Hilbert séparable de dimension infinie possède une base orthonormée
dénombrable, et concluez que deux tels espaces sont isométriquement isomorphes --- chacun d'eux est
une copie de $ \ell^2 $.
\end{enumerate}

\end{problem}

\noindent La question (c) est l'une des grandes unifications de l'analyse : l'espace des suites
de carré sommable, l'espace $ L^2 $ des fonctions de carré intégrable et l'espace des états d'une
particule quantique sont un seul et même espace de Hilbert vêtu de coordonnées différentes. Le résultat
est pour l'essentiel le théorème de Riesz--Fischer de 1907, démontré en quelques semaines par les deux
hommes dès que l'intégrale de Lebesgue eut rendu $ L^2 $ complet. Notez le contraste avec les espaces
de Banach, où aucune unicité de ce genre ne vaut --- l'espace de Hilbert est rigide à cause de son
produit scalaire.

\begin{refs}
\item Contexte : Base orthonormée --- Wikipédia (\url{https://fr.wikipedia.org/wiki/Base_orthonorm%C3%A9e})
\item Contexte : Base de Schauder --- Wikipédia (\url{https://fr.wikipedia.org/wiki/Base_de_Schauder})
\item Lecture : M. Reed et B. Simon, \textit{Methods of Modern Mathematical Physics I: Functional Analysis}, ch. II
\item Cours : MIT OCW 18.102 Introduction to Functional Analysis (en anglais) (\url{https://ocw.mit.edu/courses/18-102-introduction-to-functional-analysis-spring-2021/})
\end{refs}

\bigskip

\begin{problem}[{Les séries de Fourier comme géométrie : Parseval --- Licence}]
\textbf{FA-05.} \textbf{Problème.} On travaille dans l'espace de Hilbert $ L^2(-\pi, \pi) $ muni du produit
scalaire $ \langle f, g\rangle = \int_{-\pi}^{\pi} f(x)\,g(x)\,dx $.
\begin{enumerate}
\item[(a)] Vérifiez que le système trigonométrique
$ \big\{\, 1/\sqrt{2\pi},\ \cos(nx)/\sqrt{\pi},\ \sin(nx)/\sqrt{\pi} : n \ge 1 \,\big\} $
est orthonormé.
\item[(b)] Montrez que sa totalité équivaut à l'énoncé suivant : l'identité de Parseval vaut pour toute
$ f \in L^2(-\pi,\pi) $, c'est-à-dire que la série de Fourier de $ f $ converge vers $ f $
en norme $ L^2 $ et que la somme des carrés des coefficients de Fourier vaut $ \|f\|^2 $.
\item[(c)] En admettant la totalité (elle découle du théorème d'approximation de Weierstrass), calculez les
coefficients de Fourier de $ f(x) = x $ et déduisez-en l'évaluation d'Euler
\[ \sum_{n=1}^{\infty} \frac{1}{n^2} = \frac{\pi^2}{6} . \]
\end{enumerate}

\end{problem}

\noindent Fourier a affirmé en 1807 que des fonctions arbitraires se développent en séries
trigonométriques, et l'analyse a passé un siècle à décider de ce que cela pouvait signifier. La réponse
hilbertienne est la plus limpide : un développement de Fourier n'est rien d'autre que l'écriture d'un
vecteur dans une base orthonormée, et l'identité de Parseval (nommée d'après une observation formelle
de Marc-Antoine Parseval en 1799) est le théorème de Pythagore en dimension infinie. La question (c)
retrouve le problème de Bâle --- résolu par Euler en 1734 et traité en grande pompe sur notre page
\textit{L'art de la démonstration} --- en deux lignes de géométrie, ce qui est peut-être la meilleure
publicité qui soit pour le point de vue $ L^2 $.

\begin{refs}
\item Contexte : Égalité de Parseval --- Wikipédia (\url{https://fr.wikipedia.org/wiki/%C3%89galit%C3%A9_de_Parseval})
\item Contexte : Série de Fourier --- Wikipédia (\url{https://fr.wikipedia.org/wiki/S%C3%A9rie_de_Fourier})
\item Lecture : E. Kreyszig, \textit{Introductory Functional Analysis with Applications}, ch. 3
\item Cours : MIT OCW 18.103 Fourier Analysis (en anglais) (\url{https://ocw.mit.edu/courses/18-103-fourier-analysis-fall-2013/})
\end{refs}

\bigskip

\begin{problem}[{Opérateurs bornés et norme d'opérateur --- Licence}]
\textbf{FA-06.} \textbf{Problème.} Soient $ X, Y $ des espaces normés et $ T : X \to Y $ linéaire. On définit la
norme d'opérateur
\[ \|T\| = \sup\{\, \|Tx\| : \|x\| \le 1 \,\}, \]
et on dit que $ T $ est borné si $ \|T\| < \infty $.
\begin{enumerate}
\item[(a)] Démontrez que $ T $ est borné si et seulement si $ T $ est continu, si et seulement si
$ T $ est continu en $ 0 $.
\item[(b)] Démontrez que l'espace $ B(X, Y) $ des opérateurs bornés, muni de la norme ci-dessus, est un
espace normé, et qu'il est complet dès que $ Y $ est complet.
\item[(c)] Pour une suite bornée $ a = (a_n) $, l'opérateur de multiplication $ M_a $ sur $ \ell^2 $
envoie $ (x_n) $ sur $ (a_n x_n) $ ; démontrez que $ \|M_a\| = \sup_n |a_n| $.
\item[(d)] Pour un noyau continu $ k $ sur $ [0,1] \times [0,1] $, l'opérateur intégral
$ (Tf)(x) = \int_0^1 k(x,y)\,f(y)\,dy $ agit sur $ (C[0,1], \|\cdot\|_\infty) $ ; démontrez que
$ \|T\| = \max_{0 \le x \le 1} \int_0^1 |k(x,y)|\,dy $.
\end{enumerate}

\end{problem}

\noindent En dimension infinie la linéarité n'entraîne plus la continuité : il faut donc isoler
les opérateurs bornés ; la norme d'opérateur, implicite dans les travaux de Fredholm de 1900 sur les
équations intégrales et formalisée par Riesz et Banach, mesure l'étirement dans le pire des cas. La
question (b) est ce qui permet à la théorie spectrale de seulement commencer : $ B(X) $ est
lui-même un espace de Banach, et même une algèbre de Banach, si bien que l'on peut faire de l'analyse
--- séries entières, exponentielles, inverses via les séries de Neumann --- avec des opérateurs pour
variables.

\begin{refs}
\item Contexte : Opérateur borné --- Wikipédia (\url{https://fr.wikipedia.org/wiki/Op%C3%A9rateur_born%C3%A9})
\item Contexte : Norme d'opérateur --- Wikipédia (\url{https://fr.wikipedia.org/wiki/Norme_d'op%C3%A9rateur})
\item Lecture : E. Kreyszig, \textit{Introductory Functional Analysis with Applications}, ch. 2
\item Lecture : M. Reed et B. Simon, \textit{Methods of Modern Mathematical Physics I: Functional Analysis}, ch. III
\end{refs}

\bigskip

\begin{problem}[{Le théorème du point fixe de Banach, et pourquoi les équations différentielles ont des solutions --- Licence}]
\textbf{FA-07.} \textbf{Problème.} Le principe de contraction, et la raison pour laquelle les problèmes de Cauchy ont
des solutions uniques.
\begin{enumerate}
\item[(a)] Démontrez le théorème du point fixe de Banach : si $ (M, d) $ est un espace métrique complet non
vide et si $ T : M \to M $ vérifie $ d(Tx, Ty) \le q\,d(x, y) $ pour un $ q < 1 $ fixé,
alors $ T $ admet exactement un point fixe $ x^{*} $. Les itérées
$ x_0,\ Tx_0,\ T^2x_0, \dots $ convergent vers $ x^{*} $ quel que soit le point de départ
$ x_0 $, avec
\[ d(T^n x_0, x^{*}) \le \frac{q^n}{1-q}\, d(Tx_0, x_0) . \]
\item[(b)] Appliquez-le pour démontrer le théorème de Picard--Lindelöf (Cauchy--Lipschitz) : si $ F(x, y) $
est continue sur un rectangle autour de $ (x_0, y_0) $ et lipschitzienne en $ y $, alors le
problème de Cauchy $ y' = F(x, y) $, $ y(x_0) = y_0 $ admet exactement une solution sur un
intervalle autour de $ x_0 $.
\end{enumerate}

\end{problem}

\noindent La méthode des approximations successives d'Émile Picard (1890), affinée par Ernst
Lindelöf (1894), précède le théorème abstrait ; l'apport de Banach (1922) fut de voir qu'un seul court
argument sur les espaces métriques complets la sous-tend --- et sous-tend aussi le théorème d'inversion
locale, la convergence des itérations de type Newton et la théorie de stabilité des fractales. C'est
l'exemple le plus pur de la méthode fonctionnelle-analytique : transformer une question locale
difficile sur une équation différentielle en une question globale facile sur une application d'un
espace de fonctions.

\begin{refs}
\item Contexte : Théorème du point fixe de Banach --- Wikipédia (\url{https://fr.wikipedia.org/wiki/Th%C3%A9or%C3%A8me_du_point_fixe_de_Banach})
\item Contexte : Théorème de Cauchy-Lipschitz (Picard--Lindelöf) --- Wikipédia (\url{https://fr.wikipedia.org/wiki/Th%C3%A9or%C3%A8me_de_Cauchy-Lipschitz})
\item Lecture : E. Kreyszig, \textit{Introductory Functional Analysis with Applications}, ch. 5
\end{refs}

\bigskip

\begin{problem}[{Espaces duaux : qui sont les formes linéaires ? --- Licence}]
\textbf{FA-08.} \textbf{Problème.} Le dual $ X^{*} $ d'un espace normé $ X $ est l'espace de Banach des formes
linéaires bornées sur $ X $, muni de la norme d'opérateur.
\begin{enumerate}
\item[(a)] Démontrez que le dual de $ \ell^1 $ est isométriquement isomorphe à $ \ell^\infty $, via le
couplage $ (x, a) \mapsto \sum_n a_n x_n $.
\item[(b)] Démontrez que le dual de $ c_0 $ (suites tendant vers $ 0 $, norme sup) est isométriquement
isomorphe à $ \ell^1 $.
\item[(c)] Montrez que l'application analogue de $ \ell^1 $ dans le dual de $ \ell^\infty $ est une
isométrie mais n'est \textit{pas} surjective. (Pour les formes linéaires manquantes, voir FA-09.)
\end{enumerate}

\end{problem}

\noindent Identifier concrètement les duaux fut l'une des activités fondatrices du domaine :
F. Riesz a identifié les duaux de $ L^p $ et de $ C[0,1] $ entre 1907 et 1910, ce dernier résultat
(les formes linéaires sont des intégrations contre des mesures) étant le \og{} théorème de représentation
de Riesz \fg{} original. L'asymétrie de la question (c) est le premier signe que la dualité est profonde :
passer au dual peut agrandir un espace, la réflexivité peut échouer, et l'écart se mesure par des
objets (mesures finiment additives, limites de Banach) qu'aucune formule explicite ne produit.

\begin{refs}
\item Contexte : Dual topologique --- Wikipédia (\url{https://fr.wikipedia.org/wiki/Dual_topologique})
\item Lecture : J. B. Conway, \textit{A Course in Functional Analysis}, ch. III
\item Lecture : H. Brezis, \textit{Analyse fonctionnelle : théorie et applications}, ch. 1
\end{refs}

\bigskip

\begin{problem}[{Dimension finie : toutes les normes s'accordent --- Licence}]
\textbf{FA-19.} \textbf{Problème.} Deux normes $ \|\cdot\| $ et $ \|\cdot\|' $ sur un espace vectoriel $ X $
sont \textit{équivalentes} s'il existe des constantes $ c, C > 0 $ telles que
\[ c\,\|x\| \le \|x\|' \le C\,\|x\| \quad \text{pour tout } x \in X . \]
\begin{enumerate}
\item[(a)] Démontrez que deux normes quelconques sur $ \mathbb{R}^n $ sont équivalentes. (Comparez une
norme arbitraire à la norme euclidienne, en utilisant la compacité de la sphère unité euclidienne.)
\item[(b)] Déduisez-en que tout espace normé de dimension finie est complet, et que tout sous-espace de
dimension finie d'un espace normé est fermé.
\item[(c)] Déduisez-en que toute application linéaire d'un espace normé de dimension finie vers un espace
normé quelconque est bornée.
\item[(d)] Démontrez la dichotomie réciproque, à l'aide du lemme de Riesz (FA-12(c)) : un espace normé dont
la boule unité fermée est compacte est de dimension finie.
\end{enumerate}

\end{problem}

\noindent Ce théorème est la frontière précise entre l'algèbre linéaire et l'analyse
fonctionnelle : en deçà, la topologie est invisible --- toute norme sur $ \mathbb{R}^n $ raconte la
même histoire de convergence, et toute application linéaire est continue --- tandis qu'au-delà rien de
tout cela ne survit, comme FA-01 le montre sur $ C[0,1] $. La dichotomie de compacité de la question
(d) est due à F. Riesz (1918), et elle explique pourquoi les exemples de base du domaine sont des
espaces de fonctions : la dimension infinie, c'est exactement l'échec de Borel--Lebesgue.

\begin{refs}
\item Contexte : Norme (mathématiques) --- Wikipédia (\url{https://fr.wikipedia.org/wiki/Norme_(math%C3%A9matiques)})
\item Contexte : Lemme de Riesz --- Wikipédia (\url{https://fr.wikipedia.org/wiki/Lemme_de_Riesz})
\item Lecture : E. Kreyszig, \textit{Introductory Functional Analysis with Applications}, ch. 2
\item Cours : MIT OCW 18.102 Introduction to Functional Analysis (en anglais) (\url{https://ocw.mit.edu/courses/18-102-introduction-to-functional-analysis-spring-2021/})
\end{refs}

\bigskip

\begin{problem}[{La série de Neumann --- Licence}]
\textbf{FA-20.} \textbf{Problème.} Soit $ X $ un espace de Banach et $ T $ un opérateur borné sur $ X $ avec
$ \|T\| < 1 $.
\begin{enumerate}
\item[(a)] Démontrez que $ I - T $ est inversible dans $ B(X) $, que son inverse est la somme de la
\textit{série de Neumann}
\[ (I-T)^{-1} = \sum_{k=0}^{\infty} T^{k} , \]
et que $ \|(I-T)^{-1}\| \le (1 - \|T\|)^{-1} $.
\item[(b)] Déduisez-en que les opérateurs inversibles forment une partie ouverte de $ B(X) $, sur laquelle
l'inversion est continue.
\item[(c)] Appliquez (a) sur $ (C[0,1], \|\cdot\|_\infty) $ : si un noyau continu $ k $ vérifie
$ \max_x \int_0^1 |k(x,y)|\,dy < 1 $, alors pour toute fonction continue $ g $ l'équation
intégrale $ f(x) = g(x) + \int_0^1 k(x,y)\,f(y)\,dy $ admet exactement une solution continue
$ f $.
\item[(d)] Montrez que l'opérateur de Volterra $ (Vf)(x) = \int_0^x f(y)\,dy $ sur $ C[0,1] $ vérifie
$ \|V^n\| \le 1/n! $, et concluez que $ I - \lambda V $ est inversible pour \textit{tout} scalaire
$ \lambda $ --- l'astuce de la série porte bien au-delà de $ \|\lambda V\| < 1 $.
\end{enumerate}

\end{problem}

\noindent C'est la série géométrique avec un opérateur pour raison, nommée d'après Carl
Neumann, qui l'employa dans les années 1870 sur les équations intégrales de la théorie du potentiel ;
les approximations successives de Picard (FA-07) et les noyaux itérés de la question (c) sont la même
idée sous d'autres habits. La question (b) fait des inversibles un groupe ouvert, et c'est là que
démarre la théorie spectrale : cela force le spectre de tout opérateur borné à être un compact. La
question (d) exhibe un opérateur \textit{quasi-nilpotent} --- le spectre de $ V $ est le seul point
$ 0 $, aussi loin de la norme $ \|V\| = 1 $ que la formule du rayon spectral l'autorise.

\begin{refs}
\item Contexte : Série de Neumann --- Wikipédia (\url{https://fr.wikipedia.org/wiki/S%C3%A9rie_de_Neumann})
\item Contexte : Opérateur de Volterra --- Wikipédia (\url{https://fr.wikipedia.org/wiki/Op%C3%A9rateur_de_Volterra})
\item Lecture : E. Kreyszig, \textit{Introductory Functional Analysis with Applications}, ch. 7
\end{refs}

\bigskip

\begin{problem}[{Noyau fermé, forme linéaire bornée --- Licence, short}]
\textbf{FA-21.} \textbf{Problème.} Soit $ f $ une forme linéaire non nulle sur un espace normé
$ X $. Démontrez que $ f $ est bornée si et seulement si son noyau
$ \{\, x \in X : f(x) = 0 \,\} $ est fermé dans $ X $.
\end{problem}

\noindent La géométrie derrière l'énoncé est brutale : le noyau d'une forme linéaire est soit
fermé, soit dense dans $ X $ --- rien entre les deux. Des formes linéaires non bornées existent bel et
bien sur tout espace normé de dimension infinie, mais seulement au moyen d'une base de Hamel et de
l'axiome du choix, ce qui explique que personne n'en ait jamais écrit une explicitement.

\begin{refs}
\item Contexte : Forme linéaire --- Wikipédia (\url{https://fr.wikipedia.org/wiki/Forme_lin%C3%A9aire}), Application linéaire discontinue --- Wikipédia (en anglais) (\url{https://en.wikipedia.org/wiki/Discontinuous_linear_map})
\item Lecture : J. B. Conway, \textit{A Course in Functional Analysis}, ch. III
\end{refs}

\bigskip

\begin{problem}[{Polarisation --- Licence, short}]
\textbf{FA-22.} \textbf{Problème.} Démontrez que dans un espace préhilbertien réel le produit
scalaire se retrouve à partir de sa norme via l'identité de polarisation
$ \langle x, y\rangle = \tfrac{1}{4}\big( \|x+y\|^2 - \|x-y\|^2 \big) $, et déduisez-en que toute
application linéaire préservant la norme entre espaces préhilbertiens réels préserve automatiquement
le produit scalaire.
\end{problem}

\noindent Cette identité (avec une version à quatre termes sur $\mathbb{C}$) est la raison pour laquelle
les isométries des espaces de Hilbert sont si rigides, et elle soulève la question réciproque : quelles
normes proviennent d'un produit scalaire ? Jordan et von Neumann ont répondu en 1935 --- exactement
celles qui vérifient l'identité du parallélogramme de FA-03(a).

\begin{refs}
\item Contexte : Identité de polarisation --- Wikipédia (\url{https://fr.wikipedia.org/wiki/Identit%C3%A9_de_polarisation})
\item Lecture : J. B. Conway, \textit{A Course in Functional Analysis}, ch. I
\end{refs}

\bigskip

\begin{problem}[{Réflexivité et plongement canonique --- Licence}]
\textbf{FA-30.} \textbf{Problème.} Pour un espace normé $ X $, on définit l'\textit{application canonique}
$ J : X \to X^{**} $ dans le dual du dual par
\[ (Jx)(f) \;=\; f(x), \qquad x \in X,\ f \in X^{*} . \]
On dit que $ X $ est \textit{réflexif} si $ J $ est surjective.
\begin{enumerate}
\item[(a)] Démontrez que $ J $ est une isométrie linéaire --- l'inégalité $ \|Jx\| \le \|x\| $ est
immédiate, et la réciproque demande la forme linéaire normante produite par Hahn--Banach en FA-09(c).
Déduisez-en que tout espace normé se plonge isométriquement dans un espace de Banach, et qu'un espace
normé réflexif est complet.
\item[(b)] Démontrez que tout espace de Hilbert est réflexif (utilisez le théorème de représentation de
Riesz, FA-03(d)), et que $ \ell^p $ est réflexif pour $ 1 < p < \infty $, en admettant la
dualité $ (\ell^p)^{*} \cong \ell^q $ avec $ 1/p + 1/q = 1 $.
\item[(c)] Démontrez que $ c_0 $ et $ \ell^1 $ ne sont pas réflexifs, à l'aide des duaux calculés en
FA-08. Prenez soin de montrer que c'est l'application \textit{canonique} qui n'est pas surjective, et
non simplement qu'une application quelconque échoue.
\item[(d)] Démontrez qu'un espace de Banach $ X $ est réflexif si et seulement si $ X^{*} $ est
réflexif, et que tout sous-espace fermé d'un espace réflexif est réflexif. Déduisez-en que
$ \ell^\infty $ n'est pas réflexif.
\end{enumerate}

\end{problem}

\noindent La réflexivité est la propriété qui fait qu'un espace de dimension infinie se comporte,
en dualité, comme un espace de dimension finie, et l'insistance des questions (c) et (d) sur la
surjectivité de l'application \textit{canonique} n'est pas du pédantisme : R. C. James a construit en 1951
un espace de Banach isométriquement isomorphe à son propre bidual et pourtant non réflexif, si bien que
le type d'isomorphisme de $ X^{**} $ est le mauvais invariant. La réflexivité équivaut, par un théorème
de Kakutani, à la compacité faible de la boule unité fermée --- ce dont la méthode directe du calcul des
variations a précisément besoin, et la raison pour laquelle la convergence faible de FA-14 est bien plus
utile dans $ L^2 $ que dans $ L^1 $. La caractérisation la plus profonde est le théorème de James
(démontré pour les espaces séparables en 1957 et en général en 1964) : $ X $ est réflexif si et
seulement si toute forme linéaire bornée atteint sa norme sur la boule unité fermée --- un énoncé que tout
le monde comprend et que presque personne ne sait démontrer.

\begin{refs}
\item Contexte : Espace réflexif --- Wikipédia (\url{https://fr.wikipedia.org/wiki/Espace_r%C3%A9flexif})
\item Contexte : Théorème de James --- Wikipédia (\url{https://fr.wikipedia.org/wiki/Th%C3%A9or%C3%A8me_de_James}), Espace de James --- Wikipédia (\url{https://fr.wikipedia.org/wiki/Espace_de_James})
\item Lecture : H. Brezis, \textit{Analyse fonctionnelle : théorie et applications}, ch. 3
\item Lecture : J. B. Conway, \textit{A Course in Functional Analysis}, ch. V
\item Cours : MIT OCW 18.102 Introduction to Functional Analysis (en anglais) (\url{https://ocw.mit.edu/courses/18-102-introduction-to-functional-analysis-spring-2021/})
\end{refs}

\bigskip

\begin{problem}[{Bases de Schauder --- Licence, short}]
\textbf{FA-31.} \textbf{Problème.} On dit qu'une suite $ (e_n) $ d'un espace de Banach $ X $ est
une \textit{base de Schauder} si tout $ x \in X $ admet un unique développement $ x = \sum_n a_n e_n $
convergeant en norme. Démontrez que les vecteurs de coordonnées forment une base de Schauder de
$ c_0 $ et de $ \ell^p $ pour $ 1 \le p < \infty $ mais pas de $ \ell^\infty $, et démontrez
qu'un espace de Banach possédant une base de Schauder est nécessairement séparable.
\end{problem}

\noindent Notez ce qui est demandé : l'ordre de sommation est fixé, si bien qu'une base de
Schauder est plus faible qu'une base orthonormée (FA-04) et bien différente d'une base de Hamel.
Juliusz Schauder en a construit une pour $ C[0,1] $ en 1927, et le livre de Banach de 1932 demandait
si tout espace de Banach séparable en possédait une --- le \og{} problème de la base \fg{}, annoncé dans le Livre
écossais, où Mazur promit en 1936 une oie vivante à qui le résoudrait. Per Enflo est venu chercher l'oie
en 1972 et a publié le contre-exemple dans \textit{Acta Mathematica} en 1973 : un espace de Banach séparable
n'ayant même pas la propriété d'approximation, plus faible, et donc dépourvu de toute base.

\begin{refs}
\item Contexte : Base de Schauder --- Wikipédia (\url{https://fr.wikipedia.org/wiki/Base_de_Schauder})
\item Contexte : Propriété d'approximation --- Wikipédia (\url{https://fr.wikipedia.org/wiki/Propri%C3%A9t%C3%A9_d'approximation})
\item Article : P. Enflo, \og{} A counterexample to the approximation problem in Banach spaces \fg{} (1973), \textit{Acta Mathematica} 130, 309--317
\item Lecture : Albiac \& Kalton, \textit{Topics in Banach Space Theory}, ch. 1
\end{refs}

\bigskip


\section*{Master}

\begin{problem}[{Hahn--Banach à l'œuvre --- Master}]
\textbf{FA-09.} \textbf{Problème.} Démontrez le théorème de prolongement, puis mettez-le au travail.
\begin{enumerate}
\item[(a)] Démontrez le théorème de prolongement de Hahn--Banach : si $ p $ est une fonctionnelle
sous-linéaire sur un espace vectoriel réel $ X $ et si $ f $ est une forme linéaire sur un
sous-espace $ M $ avec $ f \le p $ sur $ M $, alors $ f $ se prolonge en une forme linéaire
sur $ X $ tout entier, encore dominée par $ p $.
\item[(b)] Déduisez-en la version \og{} espaces normés \fg{} : toute forme linéaire bornée sur un sous-espace se
prolonge à l'espace entier avec la même norme.
\item[(c)] Pour tout $ x_0 \ne 0 $ d'un espace normé $ X $, produisez $ f \in X^{*} $ avec
$ \|f\| = 1 $ et $ f(x_0) = \|x_0\| $, et concluez que $ X^{*} $ sépare les points de
$ X $.
\item[(d)] Pour un sous-espace fermé $ M $ et $ x \notin M $, produisez $ f \in X^{*} $ nulle sur
$ M $ avec $ f(x) \ne 0 $.
\item[(e)] Construisez une \textit{limite de Banach} --- une forme linéaire positive $ L $ sur
$ \ell^\infty $ invariante par le décalage
$ (x_1, x_2, \dots) \mapsto (x_2, x_3, \dots) $ et coïncidant avec la limite ordinaire sur les
suites convergentes --- et utilisez-la pour exhiber une forme linéaire sur $ \ell^\infty $ ne
provenant pas de $ \ell^1 $, achevant FA-08(c).
\end{enumerate}

\end{problem}

\noindent Démontré par Hans Hahn (1927) et, sous la forme ci-dessus, par Banach (1929), c'est le
théorème qui garantit que les espaces de dimension infinie possèdent assez de formes linéaires pour y
voir ; tout argument de dualité en analyse s'appuie sur lui. C'est aussi la première rencontre honnête du
domaine avec l'axiome du choix --- le prolongement de la question (a) existe mais en général aucune formule
ne peut le nommer, et les limites de Banach sont l'exemple canonique d'un objet dont l'existence est
démontrée et l'écriture impossible.

\begin{refs}
\item Contexte : Théorème de Hahn-Banach --- Wikipédia (\url{https://fr.wikipedia.org/wiki/Th%C3%A9or%C3%A8me_de_Hahn-Banach})
\item Lecture : H. Brezis, \textit{Analyse fonctionnelle : théorie et applications}, ch. 1
\item Lecture : J. B. Conway, \textit{A Course in Functional Analysis}, ch. III
\end{refs}

\bigskip

\begin{problem}[{Bornitude uniforme, et une série de Fourier qui diverge --- Master}]
\textbf{FA-10.} \textbf{Problème.} La catégorie de Baire verse son premier dividende : une fonction continue dont la
série de Fourier diverge.
\begin{enumerate}
\item[(a)] Démontrez le principe de bornitude uniforme (Banach--Steinhaus) : si une famille d'opérateurs
bornés d'un espace de Banach $ X $ vers un espace normé est ponctuellement bornée ---
$ \sup_T \|Tx\| < \infty $ pour chaque $ x $ --- alors elle est uniformément bornée en norme
d'opérateur.
\item[(b)] Pour $ f $ continue et $ 2\pi $-périodique, notons $ S_n(f) $ la $ n $-ième somme
partielle de sa série de Fourier évaluée en $ 0 $, de sorte que
\[ S_n(f) = \frac{1}{2\pi} \int_{-\pi}^{\pi} f(t)\,D_n(t)\,dt , \qquad
D_n(t) = \frac{\sin((n+\frac{1}{2})t)}{\sin(t/2)} \]
où $ D_n $ est le noyau de Dirichlet. Démontrez que la norme de $ S_n $ comme forme linéaire sur
$ (C(\mathbb{T}), \|\cdot\|_\infty) $ vaut
$ \frac{1}{2\pi}\int_{-\pi}^{\pi} |D_n(t)|\,dt $, et que ces nombres tendent vers l'infini comme
une constante fois $ \log n $.
\item[(c)] Concluez qu'il existe une fonction continue $ 2\pi $-périodique dont la série de Fourier diverge
en $ 0 $.
\end{enumerate}

\end{problem}

\noindent Paul du Bois-Reymond a stupéfié les analystes de 1873 en construisant une telle
fonction à la main ; le théorème de Banach--Steinhaus (1927), bâti sur le théorème de catégorie de Baire
de 1899, montre que le phénomène est générique et coûte trois lignes une fois la machinerie en place.
C'est l'illustration canonique de ce qu'achète l'\og{} analyse molle \fg{} : l'existence sans la construction.
Les constantes de la question (b) sont les constantes de Lebesgue, et leur croissance logarithmique
gouverne tout ce qui concerne la convergence ponctuelle des séries de Fourier.

\begin{refs}
\item Contexte : Théorème de Banach-Steinhaus (bornitude uniforme) --- Wikipédia (\url{https://fr.wikipedia.org/wiki/Th%C3%A9or%C3%A8me_de_Banach-Steinhaus})
\item Lecture : H. Brezis, \textit{Analyse fonctionnelle : théorie et applications}, ch. 2
\item Lecture : M. Reed et B. Simon, \textit{Methods of Modern Mathematical Physics I: Functional Analysis}, ch. III
\end{refs}

\bigskip

\begin{problem}[{Application ouverte, graphe fermé, et de la continuité gratuite --- Master}]
\textbf{FA-11.} \textbf{Problème.} Deux des grands théorèmes de Banach, et un cas de continuité offerte.
\begin{enumerate}
\item[(a)] Démontrez le théorème de l'application ouverte : un opérateur borné surjectif entre espaces de
Banach envoie les ouverts sur des ouverts ; déduisez-en qu'une bijection bornée entre espaces de
Banach a un inverse borné.
\item[(b)] Déduisez-en le théorème du graphe fermé : si $ X, Y $ sont des espaces de Banach et si
$ T : X \to Y $ est linéaire de graphe fermé --- dès que $ x_n \to x $ et $ Tx_n \to y $, on a
$ y = Tx $ --- alors $ T $ est borné.
\item[(c)] Utilisez-le pour démontrer le théorème de Hellinger--Toeplitz : un opérateur linéaire $ T $
défini sur $ H $ \textit{tout entier}, espace de Hilbert, et symétrique
($ \langle Tx, y\rangle = \langle x, Ty\rangle $ pour tous $ x, y $), est automatiquement borné.
Concluez qu'un opérateur symétrique non borné ne peut jamais être défini sur l'espace entier --- ses
ennuis de domaine, repris en FA-15, sont inévitables.
\end{enumerate}

\end{problem}

\noindent Les deux théorèmes sont de Banach (avec Schauder), publiés autour de 1929--1932, et
tous deux convertissent des hypothèses qualitatives en conclusions quantitatives via la catégorie de
Baire --- l'astuce signature de l'école polonaise. Hellinger et Toeplitz ont trouvé (c) en 1910, des
décennies avant que cela ne devienne un corollaire de deux lignes ; von Neumann y a plus tard reconnu la
raison pour laquelle les observables quantiques comme l'énergie et l'impulsion, qui sont non bornées,
doivent vivre sur des domaines denses et propres. La continuité, dans un espace de Banach, est parfois
gratuite.

\begin{refs}
\item Contexte : Théorème de Banach-Schauder (application ouverte) --- Wikipédia (\url{https://fr.wikipedia.org/wiki/Th%C3%A9or%C3%A8me_de_Banach-Schauder})
\item Contexte : Théorème du graphe fermé --- Wikipédia (\url{https://fr.wikipedia.org/wiki/Th%C3%A9or%C3%A8me_du_graphe_ferm%C3%A9}), Théorème de Hellinger--Toeplitz --- Wikipédia (en anglais) (\url{https://en.wikipedia.org/wiki/Hellinger%E2%80%93Toeplitz_theorem})
\item Lecture : H. Brezis, \textit{Analyse fonctionnelle : théorie et applications}, ch. 2
\end{refs}

\bigskip

\begin{problem}[{Opérateurs compacts --- Master}]
\textbf{FA-12.} \textbf{Problème.} Un opérateur $ T $ sur un espace de Hilbert $ H $ est
\textit{compact} s'il envoie la boule unité fermée sur un ensemble d'adhérence compacte.
\begin{enumerate}
\item[(a)] Démontrez que les opérateurs compacts forment un idéal bilatère $ K(H) $ fermé en norme dans
$ B(H) $.
\item[(b)] Démontrez que toute limite en norme d'opérateurs de rang fini est compacte, et que sur un espace
de Hilbert la réciproque vaut également.
\item[(c)] Démontrez le lemme de Riesz --- pour tout sous-espace fermé propre $ M $ d'un espace normé
$ X $ et tout $ \theta < 1 $ il existe un vecteur unitaire $ x $ avec
$ \mathrm{dist}(x, M) \ge \theta $ --- et déduisez-en que l'opérateur identité d'un espace normé de
dimension infinie n'est jamais compact.
\item[(d)] Démontrez qu'un opérateur intégral sur $ L^2[0,1] $ de noyau continu $ k(x,y) $ est
compact.
\end{enumerate}

\end{problem}

\noindent Les opérateurs compacts sont les matrices infinies qui se comportent comme des
matrices finies, et c'est là que la théorie spectrale a d'abord fonctionné : la théorie des équations
intégrales de Fredholm (1900) et les articles de Hilbert de 1904--1906 sont, en langage moderne, l'étude
de $ K(H) $. La propriété d'idéal de la question (a) est le germe de toute une discipline --- l'algèbre
de Calkin $ B(H)/K(H) $, la théorie de l'indice de Fredholm, et finalement la géométrie non commutative
de FA-16 et FA-18. La question (c) explique la pathologie fondamentale de la dimension infinie : les
boules fermées ne sont jamais compactes, si bien que la compacité doit être importée par l'opérateur.

\begin{refs}
\item Contexte : Opérateur compact --- Wikipédia (\url{https://fr.wikipedia.org/wiki/Op%C3%A9rateur_compact})
\item Contexte : Lemme de Riesz --- Wikipédia (\url{https://fr.wikipedia.org/wiki/Lemme_de_Riesz}), Alternative de Fredholm --- Wikipédia (\url{https://fr.wikipedia.org/wiki/Alternative_de_Fredholm})
\item Lecture : E. Kreyszig, \textit{Introductory Functional Analysis with Applications}, ch. 8
\item Lecture : J. B. Conway, \textit{A Course in Functional Analysis}, ch. II
\end{refs}

\bigskip

\begin{problem}[{Le théorème spectral pour les opérateurs compacts autoadjoints --- Master}]
\textbf{FA-13.} \textbf{Problème.} Soit $ T $ un opérateur compact autoadjoint sur un espace de Hilbert
$ H $.
\begin{enumerate}
\item[(a)] Démontrez que $ \|T\| = \sup\{\, |\langle Tx, x\rangle| : \|x\| = 1 \,\} $ et que cette borne
supérieure est atteinte, de sorte que $ \|T\| $ ou $ -\|T\| $ est valeur propre de
$ T $.
\item[(b)] Démontrez que les vecteurs propres associés à des valeurs propres distinctes sont orthogonaux,
que chaque valeur propre non nulle a un sous-espace propre de dimension finie, et que les valeurs
propres non nulles forment un ensemble fini ou une suite tendant vers $ 0 $.
\item[(c)] Démontrez le théorème spectral : $ H $ possède une base orthonormée de vecteurs propres de
$ T $, et
\[ T = \sum_n \lambda_n \langle \cdot\,, e_n\rangle\, e_n \]
avec convergence en norme d'opérateur.
\item[(d)] Expliquez en une page comment cette diagonalisation résout le problème de la corde vibrante : les
fonctions propres de l'opérateur intégral pertinent sont les modes $ \sin(nx) $, et le développement
selon ces modes est la série de Fourier de FA-05.
\end{enumerate}

\end{problem}

\noindent C'est le couronnement de la théorie classique, démontré par Hilbert (1906) pour les
équations intégrales et mis sous sa forme moderne par Riesz et Schmidt. C'est l'analogue exact, en
dimension infinie, de la diagonalisation d'une matrice symétrique, et sa signification physique --- tout
système symétrique \og{} raisonnable \fg{} se décompose en modes indépendants de fréquences réelles --- explique
que la théorie de Sturm--Liouville, les états liés quantiques et l'analyse en composantes principales
finissent par faire le même calcul. Tout FA-15 et FA-16 commence par se demander ce qui survit lorsqu'on
abandonne la compacité.

\begin{refs}
\item Contexte : Opérateur compact sur un espace de Hilbert --- Wikipédia (en anglais) (\url{https://en.wikipedia.org/wiki/Compact_operator_on_Hilbert_space})
\item Contexte : Théorème spectral --- Wikipédia (\url{https://fr.wikipedia.org/wiki/Th%C3%A9or%C3%A8me_spectral})
\item Lecture : H. Brezis, \textit{Analyse fonctionnelle : théorie et applications}, ch. 6
\item Cours : MIT OCW 18.102 Introduction to Functional Analysis (en anglais) (\url{https://ocw.mit.edu/courses/18-102-introduction-to-functional-analysis-spring-2021/})
\end{refs}

\bigskip

\begin{problem}[{Convergence faible --- Master}]
\textbf{FA-14.} \textbf{Problème.} Une suite $ (x_n) $ d'un espace de Hilbert $ H $
\textit{converge faiblement} vers $ x $ si $ \langle x_n, y\rangle \to \langle x, y\rangle $ pour
tout $ y \in H $.
\begin{enumerate}
\item[(a)] Montrez que toute suite orthonormée converge faiblement vers $ 0 $ mais ne converge pas en
norme --- la convergence faible et la convergence en norme diffèrent donc réellement.
\item[(b)] Démontrez que les limites faibles n'augmentent pas la norme : si $ x_n \to x $ faiblement alors
$ \|x\| \le \liminf_n \|x_n\| $, et si de plus $ \|x_n\| \to \|x\| $ alors $ x_n \to x $ en
norme (propriété de Radon--Riesz).
\item[(c)] Démontrez que toute suite bornée d'un espace de Hilbert séparable admet une sous-suite faiblement
convergente.
\item[(d)] Montrez par un exemple que la limite faible de vecteurs unitaires peut être $ 0 $, et expliquez
pourquoi (c) est le bon substitut du théorème de Bolzano--Weierstrass, dont FA-12(c) a montré l'échec
en norme.
\end{enumerate}

\end{problem}

\noindent La convergence faible a été distillée par Hilbert et Riesz à partir des arguments de
\og{} principe de choix \fg{} du calcul des variations : les suites minimisantes convergent rarement, mais elles
convergent faiblement, et la semi-continuité inférieure de (b) est exactement ce dont la méthode directe
du calcul des variations a besoin. La question (c) est le cas hilbertien séparable du théorème de
Banach--Alaoglu (1940) et de la caractérisation de la réflexivité par Kakutani ; en EDP on s'en sert tous
les jours, sous le nom d'\og{} extraction d'une sous-suite faiblement convergente \fg{}, pour fabriquer des
solutions comme limites faibles.

\begin{refs}
\item Contexte : Convergence faible (espace de Hilbert) --- Wikipédia (\url{https://fr.wikipedia.org/wiki/Convergence_faible_(espace_de_Hilbert)})
\item Contexte : Théorème de Banach--Alaoglu --- Wikipédia (\url{https://fr.wikipedia.org/wiki/Th%C3%A9or%C3%A8me_de_Banach-Alaoglu-Bourbaki})
\item Lecture : H. Brezis, \textit{Analyse fonctionnelle : théorie et applications}, ch. 3
\end{refs}

\bigskip

\begin{problem}[{Opérateurs non bornés et la forme de la mécanique quantique --- Master}]
\textbf{FA-15.} \textbf{Problème.} La relation de commutation de Heisenberg ne peut pas vivre dans $ B(H) $ ; ce
problème localise là où elle vit.
\begin{enumerate}
\item[(a)] Démontrez le théorème de Wielandt--Wintner : il n'existe pas d'opérateurs bornés $ A, B $ sur un
espace normé avec $ AB - BA = I $.
\item[(b)] Sur $ L^2(\mathbb{R}) $, définissez l'opérateur de position $ (Qf)(x) = x f(x) $ sur son
domaine maximal $ D(Q) = \{ f \in L^2(\mathbb{R}) : x f(x) \in L^2(\mathbb{R}) \} $, et l'opérateur
d'impulsion $ P = -i\,\frac{d}{dx} $ sur l'espace de Schwartz des fonctions lisses à décroissance
rapide. Vérifiez que $ PQ - QP = -iI $ sur l'espace de Schwartz, et concluez de (a) que $ P $ et
$ Q $ ne peuvent pas être tous deux bornés.
\item[(c)] Démontrez que $ Q $, muni du domaine ci-dessus, est autoadjoint : $ Q = Q^{*} $, où l'adjoint
d'un opérateur à domaine dense est défini par $ \langle Tx, y\rangle = \langle x, T^{*}y\rangle $
sur le domaine naturel.
\item[(d)] Montrez que le spectre de $ Q $ est $\mathbb{R}$ tout entier bien que $ Q $ n'ait aucun vecteur propre
dans $ L^2(\mathbb{R}) $ --- le spectre est purement continu.
\end{enumerate}

\end{problem}

\noindent La relation de commutation de 1925 entre position et impulsion, chez Heisenberg, est
par (a) impossible pour des opérateurs bornés --- la mécanique quantique \textit{impose} donc la théorie des
opérateurs non bornés à domaine dense, que von Neumann a bâtie en 1929--1932 exactement à cette fin, en
traçant la ligne cruciale entre symétrique et autoadjoint. La question (d) montre le spectre débordant la
valeur propre : une mesure de position peut donner n'importe quel réel, dont aucun ne correspond à un
état normalisable. Le théorème de Stone, qui identifie les opérateurs autoadjoints aux générateurs
d'évolutions unitaires en temps, complète ce tableau et constitue le problème suivant naturel à
chercher.

\begin{refs}
\item Contexte : Opérateur non borné --- Wikipédia (\url{https://fr.wikipedia.org/wiki/Op%C3%A9rateur_non_born%C3%A9}), Opérateur autoadjoint --- Wikipédia (\url{https://fr.wikipedia.org/wiki/Endomorphisme_autoadjoint})
\item Contexte : Théorème de Stone sur les groupes unitaires à un paramètre --- Wikipédia (en anglais) (\url{https://en.wikipedia.org/wiki/Stone%27s_theorem_on_one-parameter_unitary_groups})
\item Lecture : M. Reed et B. Simon, \textit{Methods of Modern Mathematical Physics I: Functional Analysis}, ch. VIII
\item Lecture : E. Kreyszig, \textit{Introductory Functional Analysis with Applications}, ch. 10--11
\end{refs}

\bigskip

\begin{problem}[{L'identité C* : là où l'algèbre rencontre l'analyse --- Master}]
\textbf{FA-16.} \textbf{Problème.} Une \textit{C*-algèbre} est une algèbre de Banach $ A $ munie d'une involution
$ a \mapsto a^{*} $ vérifiant l'identité C*
\[ \|a^{*}a\| = \|a\|^2 . \]
\begin{enumerate}
\item[(a)] Vérifiez que $ B(H) $ muni de l'adjonction, et $ C(X) $ (fonctions complexes continues sur un
espace compact séparé, norme sup, conjugaison ponctuelle) sont des C*-algèbres.
\item[(b)] Démontrez que dans toute C*-algèbre $ \|a^{*}\| = \|a\| $.
\item[(c)] Démontrez que pour $ a $ autoadjoint on a $ \|a^2\| = \|a\|^2 $, et déduisez-en, à l'aide de
la formule du rayon spectral de Gelfand $ r(a) = \lim_n \|a^n\|^{1/n} $, que $ \|a\| = r(a) $ :
pour les éléments autoadjoints la norme est déterminée par le spectre.
\item[(d)] Concluez qu'une *-algèbre admet au plus une norme qui en fasse une C*-algèbre --- l'analyse est
entièrement encodée dans l'algèbre.
\end{enumerate}

\end{problem}

\noindent Gelfand et Naimark ont introduit ces axiomes en 1943 et démontré les deux théorèmes de
structure qui les justifient : toute C*-algèbre commutative est $ C(X) $ pour un espace compact
canonique $ X $, et toute C*-algèbre se plonge dans un $ B(H) $. La question (d) est le cœur
philosophique --- la norme, donc la topologie, donc l'\og{} espace \fg{}, se retrouve à partir de l'algèbre pure ---
et c'est pourquoi Connes a pu proposer les C*-algèbres et les algèbres de von Neumann comme \og{} topologie
non commutative \fg{} et \og{} théorie de la mesure non commutative \fg{}. Ce problème est la porte d'entrée vers
FA-18 et vers une large part de la théorie des opérateurs contemporaine.

\begin{refs}
\item Contexte : C*-algèbre --- Wikipédia (\url{https://fr.wikipedia.org/wiki/C*-alg%C3%A8bre})
\item Contexte : Théorème de Gelfand--Naimark --- Wikipédia (en anglais) (\url{https://en.wikipedia.org/wiki/Gelfand%E2%80%93Naimark_theorem}), Spectre d'un opérateur linéaire --- Wikipédia (\url{https://fr.wikipedia.org/wiki/Spectre_d'un_op%C3%A9rateur_lin%C3%A9aire})
\item Lecture : J. B. Conway, \textit{A Course in Functional Analysis}, ch. VIII
\end{refs}

\bigskip

\begin{problem}[{L'adjoint --- Master}]
\textbf{FA-23.} \textbf{Problème.} Soit $ H $ un espace de Hilbert complexe et $ T \in B(H) $. Le théorème de
représentation de Riesz (FA-03(d)) produit un unique opérateur borné $ T^{*} $, l'\textit{adjoint} de
$ T $, vérifiant
\[ \langle Tx, y\rangle = \langle x, T^{*}y\rangle \quad \text{pour tous } x, y \in H . \]
\begin{enumerate}
\item[(a)] Menez cette construction, et démontrez que $ \|T^{*}\| = \|T\| $ et
$ \|T^{*}T\| = \|T\|^{2} $.
\item[(b)] Démontrez la dualité du noyau et de l'image :
$ \ker T^{*} = (\operatorname{ran} T)^{\perp} $ et
$ \overline{\operatorname{ran} T} = (\ker T^{*})^{\perp} $ ; déduisez-en que $ T $ est d'image
dense si et seulement si $ T^{*} $ est injectif.
\item[(c)] Calculez des adjoints : l'opérateur de multiplication $ M_a $ de FA-06(c) vérifie
$ M_a^{*} = M_{\bar a} $ ; le décalage unilatéral
$ S(x_1, x_2, \dots) = (0, x_1, x_2, \dots) $ sur $ \ell^2 $ a pour adjoint le décalage arrière
$ S^{*}(x_1, x_2, \dots) = (x_2, x_3, \dots) $ ; vérifiez que
$ S^{*}S = I \ne S S^{*} $.
\item[(d)] On dit que $ T $ est \textit{autoadjoint} si $ T = T^{*} $, \textit{normal} si
$ TT^{*} = T^{*}T $, et \textit{unitaire} si $ TT^{*} = T^{*}T = I $. Démontrez que $ T $ est
normal si et seulement si $ \|Tx\| = \|T^{*}x\| $ pour tout $ x $, et que $ T $ est unitaire si
et seulement si c'est une isométrie surjective.
\end{enumerate}

\end{problem}

\noindent L'adjoint, c'est la transposée conjuguée libérée des coordonnées : il est entré en
analyse comme le \og{} noyau transposé \fg{} $ k(y,x) $ des équations intégrales de Fredholm, et la définition
abstraite de von Neumann (1929--30) est ce qui a rendu l'autoadjonction --- la propriété dont la mécanique
quantique a besoin (FA-15) --- exprimable sans matrices. L'identité $ \|T^{*}T\| = \|T\|^{2} $ démontrée
en (a) est l'identité C* : FA-16 prend exactement cette équation pour axiome et en fait repousser tout le
jeu de l'algèbre et de l'analyse dans $ B(H) $. Les opérateurs qui échouent à $ SS^{*} = S^{*}S $,
comme le décalage de (c), sont là où la théorie cesse de ressembler à l'algèbre linéaire ; FA-24 étudie
cet exemple pour de bon.

\begin{refs}
\item Contexte : Opérateur adjoint --- Wikipédia (\url{https://fr.wikipedia.org/wiki/Op%C3%A9rateur_adjoint})
\item Lecture : E. Kreyszig, \textit{Introductory Functional Analysis with Applications}, ch. 3
\item Lecture : J. B. Conway, \textit{A Course in Functional Analysis}, ch. II
\end{refs}

\bigskip

\begin{problem}[{Le décalage et son spectre --- Master}]
\textbf{FA-24.} \textbf{Problème.} Le \textit{spectre} $ \sigma(T) $ de $ T \in B(H) $ est l'ensemble des
$ \lambda \in \mathbb{C} $ pour lesquels $ T - \lambda I $ n'a pas d'inverse borné. On travaille
sur $ \ell^2 $ avec le décalage unilatéral $ S $ et le décalage arrière $ S^{*} $ de
FA-23(c).
\begin{enumerate}
\item[(a)] Montrez que $ S $ est une isométrie non unitaire, et que $ S $ n'a aucune valeur
propre.
\item[(b)] Montrez que tout $ \lambda $ avec $ |\lambda| < 1 $ est valeur propre de $ S^{*} $, de
vecteur propre $ (1, \lambda, \lambda^2, \dots) $.
\item[(c)] À l'aide de (b), de la fermeture du spectre et de la borne
$ \sigma(T) \subset \{\, \lambda : |\lambda| \le \|T\| \,\} $ fournie par la série de Neumann
(FA-20), démontrez que
\[ \sigma(S) = \sigma(S^{*}) = \{\, \lambda \in \mathbb{C} : |\lambda| \le 1 \,\} . \]
\item[(d)] Prenez maintenant le décalage bilatéral $ U(x_n)_{n \in \mathbb{Z}} = (x_{n-1})_{n \in \mathbb{Z}} $
sur $ \ell^2(\mathbb{Z}) $. Démontrez que $ U $ est unitaire et que $ \sigma(U) $ est exactement
le cercle unité. (Pour $ |\lambda| = 1 $, construisez des vecteurs propres approchés à partir de
longues troncatures de $ (\dots, \lambda^{-1}, 1, \lambda, \lambda^{2}, \dots) $.)
\end{enumerate}

\end{problem}

\noindent Halmos appelait le décalage unilatéral l'opérateur le plus important de l'espace de
Hilbert : c'est le plus simple opérateur sans valeur propre mais avec un disque entier de spectre, et par
la décomposition de Wold il est la brique de toute isométrie --- chacune est une partie unitaire plus des
copies de $ S $. Ses sous-espaces invariants ont été classifiés par Beurling à la fin des années 1940
en termes de fonctions intérieures de l'espace de Hardy $ H^2 $, pont précoce entre théorie des
opérateurs et analyse complexe et raison pour laquelle les opérateurs de type décalage sont au cœur des
travaux sur le problème du sous-espace invariant (FA-17). Le contraste entre (c) et (d) --- même formule,
espace différent, spectre différent --- est l'avertissement le plus net que les spectres sont globaux et
non formulaires.

\begin{refs}
\item Contexte : Opérateur de décalage --- Wikipédia (\url{https://fr.wikipedia.org/wiki/Op%C3%A9rateur_de_d%C3%A9calage})
\item Contexte : Spectre d'un opérateur linéaire --- Wikipédia (\url{https://fr.wikipedia.org/wiki/Spectre_d'un_op%C3%A9rateur_lin%C3%A9aire})
\item Lecture : P. R. Halmos, \textit{A Hilbert Space Problem Book}
\end{refs}

\bigskip

\begin{problem}[{L'alternative de Fredholm --- Master}]
\textbf{FA-25.} \textbf{Problème.} Soit $ H $ un espace de Hilbert et $ K $ un opérateur compact sur $ H $
(FA-12).
\begin{enumerate}
\item[(a)] Démontrez que $ \ker(I - K) $ est de dimension finie et que $ \operatorname{ran}(I - K) $ est
fermé.
\item[(b)] Démontrez l'\textit{alternative de Fredholm} : $ I - K $ est injectif si et seulement s'il est
surjectif. Ainsi, ou bien l'équation homogène $ x - Kx = 0 $ a une solution non nulle, ou bien
$ x - Kx = y $ est uniquement soluble pour tout $ y $.
\item[(c)] Démontrez que $ \dim \ker(I - K) = \dim \ker(I - K^{*}) $, de sorte que le nombre de conditions
de solubilité portant sur $ y $ égale le nombre de degrés de liberté du problème homogène.
\item[(d)] Déduisez-en le tableau spectral qui complète FA-13 : tout point non nul de $ \sigma(K) $ est une
valeur propre de multiplicité finie. Traduisez ensuite (b)--(c) en la dichotomie classique de
solubilité de l'équation intégrale
$ f(x) - \lambda \int_0^1 k(x,y)\,f(y)\,dy = g(x) $ à noyau continu $ k $.
\end{enumerate}

\end{problem}

\noindent L'article d'Ivar Fredholm paru en 1903 dans \textit{Acta Mathematica} a établi cette
dichotomie pour les équations intégrales en prenant au sérieux le déterminant d'une matrice infinie, et
il a fait détonation : Hilbert a passé les années suivantes à bâtir ce qui allait devenir la théorie
spectrale, et Riesz a reformulé l'alternative pour les opérateurs compacts abstraits en 1918 --- la forme
démontrée ici. L'alternative est l'écho, en dimension infinie, du \og{} injectif si et seulement si
surjectif \fg{} de l'algèbre linéaire, restauré par la compacité après que FA-24 en eut montré l'échec cuisant
en général ($ S $ est injectif et jamais surjectif). En EDP c'est la machine standard pour le problème
de Dirichlet : l'existence pour l'équation intégrale de bord suit dès l'unicité vérifiée.

\begin{refs}
\item Contexte : Alternative de Fredholm --- Wikipédia (\url{https://fr.wikipedia.org/wiki/Alternative_de_Fredholm})
\item Lecture : H. Brezis, \textit{Analyse fonctionnelle : théorie et applications}, ch. 6
\item Article : I. Fredholm, \og{} Sur une classe d'équations fonctionnelles \fg{} (1903), \textit{Acta Mathematica} 27
\end{refs}

\bigskip

\begin{problem}[{Projecteurs et sous-espaces --- Master, short}]
\textbf{FA-26.} \textbf{Problème.} Soit $ P $ un opérateur borné sur un espace de Hilbert avec
$ P^2 = P = P^{*} $. Démontrez que $ P $ est la projection orthogonale sur le sous-espace fermé
$ \operatorname{ran} P $ au sens de FA-03(c), et que tout sous-espace fermé provient de cette manière
d'un unique tel $ P $.
\end{problem}

\noindent Ce dictionnaire --- sous-espaces fermés d'un côté, idempotents autoadjoints de l'autre ---
est celui de von Neumann, et il porte le bâtiment : le théorème spectral écrit tout opérateur autoadjoint
comme une intégrale de projecteurs, et la mécanique quantique lit $ \langle Px, x\rangle $ comme la
probabilité de l'événement $ P $. Abandonner la condition $ P = P^{*} $ laisse les projecteurs
obliques, dont les normes peuvent être arbitrairement grandes.

\begin{refs}
\item Contexte : Projecteur (algèbre linéaire) --- Wikipédia (\url{https://fr.wikipedia.org/wiki/Projecteur_(math%C3%A9matiques)})
\item Lecture : J. B. Conway, \textit{A Course in Functional Analysis}, ch. II
\end{refs}

\bigskip

\begin{problem}[{Le principe de sélection de Helly : Banach--Alaoglu en version séquentielle --- Master, short}]
\textbf{FA-27.} \textbf{Problème.} Soit $ X $ un espace normé séparable. Démontrez que toute suite
bornée $ (f_n) $ du dual $ X^{*} $ admet une sous-suite faiblement-* convergente : il existe des
indices $ n_1 < n_2 < \cdots $ et $ f \in X^{*} $ avec $ f_{n_k}(x) \to f(x) $ pour tout
$ x \in X $.
\end{problem}

\noindent Cette compacité par argument diagonal, trouvée par Helly en 1912 pour les formes
linéaires sur $ C[a,b] $, est le germe séquentiel du théorème de Banach--Alaoglu (1940), dont la forme
complète --- la boule unité fermée de tout dual est faiblement-* compacte --- exige le théorème de Tykhonov.
C'est le pain quotidien des EDP et des probabilités : \og{} passer à une limite faible-* \fg{} fabrique des
mesures, des solutions faibles et des limites généralisées chaque fois que la convergence explicite est
hors de portée.

\begin{refs}
\item Contexte : Théorème de Banach--Alaoglu --- Wikipédia (\url{https://fr.wikipedia.org/wiki/Th%C3%A9or%C3%A8me_de_Banach-Alaoglu-Bourbaki})
\item Lecture : H. Brezis, \textit{Analyse fonctionnelle : théorie et applications}, ch. 3
\end{refs}

\bigskip

\begin{problem}[{La formule du rayon spectral --- Master, short}]
\textbf{FA-32.} \textbf{Problème.} Soit $ A $ une algèbre de Banach complexe unifère et
$ a \in A $, et posons
$ \sigma(a) = \{\, \lambda \in \mathbb{C} : \lambda 1 - a \text{ n'est pas inversible} \,\} $.
Démontrez que $ \sigma(a) $ est un compact non vide de $ \mathbb{C} $, et que son rayon
$ r(a) = \sup\{\, |\lambda| : \lambda \in \sigma(a) \,\} $ est donné par la formule de
Beurling--Gelfand
\[ r(a) \;=\; \lim_{n \to \infty} \|a^n\|^{1/n} \;=\; \inf_{n \ge 1} \|a^n\|^{1/n} . \]
\end{problem}

\noindent La compacité vient de la série de Neumann de FA-20, mais la non-vacuité est le point où
l'analyse complexe entre dans l'analyse fonctionnelle : la résolvante
$ \lambda \mapsto (\lambda 1 - a)^{-1} $ est une fonction holomorphe à valeurs dans un espace de Banach
qui s'annule à l'infini, si bien qu'un spectre vide contredirait le théorème de Liouville. La formule, due
à Gelfand en 1941 et anticipée par Beurling en 1938, est remarquable en ce qu'elle égale une quantité
purement spectrale à une quantité purement métrique ; l'opérateur de Volterra de FA-20(d), avec
$ \|V\| = 1 $ et $ r(V) = 0 $, montre à quel point les deux membres peuvent s'écarter pour un $ n $
donné. FA-16(c) a utilisé cette formule pour démontrer qu'une C*-algèbre n'admet qu'une seule norme
possible --- ce problème en fournit l'outil.

\begin{refs}
\item Contexte : Rayon spectral --- Wikipédia (\url{https://fr.wikipedia.org/wiki/Rayon_spectral})
\item Contexte : Algèbre de Banach --- Wikipédia (\url{https://fr.wikipedia.org/wiki/Alg%C3%A8bre_de_Banach})
\item Lecture : W. Rudin, \textit{Functional Analysis}, ch. 10
\end{refs}

\bigskip

\begin{problem}[{La théorie de Gelfand et le lemme de Wiener --- Master}]
\textbf{FA-33.} \textbf{Problème.} Soit $ A $ une algèbre de Banach complexe commutative unifère. Un
\textit{caractère} de $ A $ est une forme linéaire multiplicative non nulle
$ \varphi : A \to \mathbb{C} $ ; on note $ \Delta(A) $ l'ensemble des caractères.
\begin{enumerate}
\item[(a)] Déduisez le théorème de Gelfand--Mazur de la non-vacuité du spectre (FA-32) : une algèbre de Banach
complexe unifère dont tout élément non nul est inversible est isométriquement isomorphe à
$ \mathbb{C} $.
\item[(b)] Démontrez que tout caractère est automatiquement continu avec $ \|\varphi\| = 1 $, et que
$ \varphi \mapsto \ker\varphi $ est une bijection de $ \Delta(A) $ sur l'ensemble des idéaux
maximaux de $ A $. (Les idéaux maximaux sont fermés ; le quotient par l'un d'eux est une algèbre à
division de Banach ; appliquez (a).)
\item[(c)] Munissez $ \Delta(A) $ de la topologie faible-* héritée de $ A^{*} $ et démontrez qu'elle est
compacte séparée (Banach--Alaoglu, FA-27). Démontrez que la \textit{transformation de Gelfand}
\[ \hat{a}(\varphi) = \varphi(a), \qquad a \mapsto \hat{a} \in C(\Delta(A)), \]
est un morphisme d'algèbres décroissant la norme, avec $ \hat{a}(\Delta(A)) = \sigma(a) $ et
$ \|\hat{a}\|_\infty = r(a) $, et identifiez son noyau.
\item[(d)] Faites tourner la machine. Soit $ A $ l'algèbre de Wiener des fonctions
$ f(\theta) = \sum_{n \in \mathbb{Z}} c_n e^{in\theta} $ avec $ \sum_n |c_n| < \infty $, qui est
$ \ell^1(\mathbb{Z}) $ muni du produit de convolution. Démontrez que les caractères de $ A $ sont
exactement les évaluations $ f \mapsto f(\theta_0) $, de sorte que $ \Delta(A) $ est le cercle, et
déduisez-en le lemme de Wiener : si $ f $ a une série de Fourier absolument convergente et ne
s'annule jamais, alors $ 1/f $ a également une série de Fourier absolument convergente.
\end{enumerate}

\end{problem}

\noindent Israel Gelfand a créé cette théorie dans son article de 1941 sur les anneaux normés, et
la question (d) en fut la carte de visite : Wiener avait démontré le lemme en 1932 par un argument long et
délicat au sein de ses travaux sur les théorèmes taubériens, et la démonstration de Gelfand tient en
quatre lignes une fois que l'on connaît les idéaux maximaux. La leçon se généralise --- une algèbre de
Banach commutative est un espace de fonctions sur son propre espace d'idéaux maximaux, l'inversibilité
signifie la non-annulation, et l'analyse se convertit en topologie. Ajouter l'identité C* de FA-16 rend la
transformation isométrique et surjective, ce qui est le théorème de Gelfand--Naimark commutatif
$ A \cong C(X) $ et la raison pour laquelle un espace compact séparé et son algèbre de fonctions sont
interchangeables --- le point de départ de la géométrie non commutative.

\begin{refs}
\item Contexte : Transformation de Gelfand --- Wikipédia (en anglais) (\url{https://en.wikipedia.org/wiki/Gelfand_representation})
\item Contexte : Théorème de Gelfand-Mazur --- Wikipédia (\url{https://fr.wikipedia.org/wiki/Th%C3%A9or%C3%A8me_de_Gelfand-Mazur}), Lemme de Wiener --- Wikipédia (en anglais) (\url{https://en.wikipedia.org/wiki/Wiener%27s_lemma})
\item Lecture : W. Rudin, \textit{Functional Analysis}, ch. 11
\item Lecture : J. B. Conway, \textit{A Course in Functional Analysis}, ch. VII
\end{refs}

\bigskip

\begin{problem}[{L'espace de Schwartz et les distributions tempérées --- Master}]
\textbf{FA-34.} \textbf{Problème.} Soit $ \mathcal{S}(\mathbb{R}) $ l'espace des fonctions lisses dont toutes les
dérivées décroissent plus vite que toute puissance, topologisé par les semi-normes
$ p_{k,m}(f) = \sup_{x} |x^k f^{(m)}(x)| $, et utilisez la transformée de Fourier dans la
normalisation
\[ \hat{f}(\xi) \;=\; \int_{\mathbb{R}} f(x)\, e^{-2\pi i x\xi}\, dx . \]
\begin{enumerate}
\item[(a)] Démontrez que $ \mathcal{S}(\mathbb{R}) $ est un espace de Fréchet --- sa topologie provient d'une
distance invariante par translation pour laquelle il est complet --- et qu'il n'est pas normable. Montrez
que $ C_c^{\infty}(\mathbb{R}) $ y est dense.
\item[(b)] Démontrez que la transformée de Fourier envoie $ \mathcal{S} $ continûment dans
$ \mathcal{S} $, échangeant $ f' $ et $ 2\pi i \xi \hat{f} $, ainsi que $ xf $ et
$ -\hat{f}'/(2\pi i) $ ; calculez la transformée de la gaussienne $ e^{-\pi x^2} $ ; et démontrez
le théorème d'inversion, sous la forme suivante : la transformée de Fourier est une bijection de
$ \mathcal{S} $ sur lui-même dont la puissance quatrième est l'identité.
\item[(c)] Démontrez le théorème de Plancherel : $ \|\hat{f}\|_{L^2} = \|f\|_{L^2} $ pour
$ f \in \mathcal{S} $, et déduisez-en que la transformée de Fourier se prolonge de manière unique en
un opérateur unitaire sur $ L^2(\mathbb{R}) $.
\item[(d)] Une \textit{distribution tempérée} est une forme linéaire continue sur $ \mathcal{S} $. Démontrez
que l'évaluation $ \delta(f) = f(0) $ en est une, et qu'aucune fonction localement intégrable
$ g $ ne vérifie $ \int g f = f(0) $ pour toute $ f \in \mathcal{S} $. Démontrez que la dérivée
au sens des distributions de la fonction de Heaviside est $ \delta $, prolongez la transformée de
Fourier aux distributions tempérées par dualité, et calculez les transformées de $ \delta $ et de la
fonction constante 1.
\end{enumerate}

\end{problem}

\noindent Physiciens et ingénieurs dérivaient des fonctions discontinues et utilisaient le $\delta$ de
Dirac depuis deux décennies lorsque Laurent Schwartz, dans sa \textit{Théorie des distributions} (1950--51),
donna un foyer à ces opérations : une distribution est une forme linéaire, la dérivation est transposée
sur la fonction test, et toute distribution est infiniment dérivable par décret. Sobolev avait eu l'idée
essentielle en 1936 dans ses travaux sur l'équation des ondes, et Schwartz a reçu la médaille Fields en
1950 pour la théorie générale. Notez ce qu'achète l'espace $ \mathcal{S} $ à la question (c) : la
transformée de Fourier est une symétrie parfaite de $ \mathcal{S} $, si bien qu'elle se prolonge par
dualité à $ \mathcal{S}' $ et par densité à $ L^2 $, et c'est seulement alors que les transformées de
$ \delta $, des polynômes et de $ e^{2\pi i a x} $ prennent un sens. C'est aussi le moment où la
discipline quitte les espaces normés : $ \mathcal{S} $ est métrisable mais non normable, et c'est bien
pourquoi les espaces localement convexes existent.

\begin{refs}
\item Contexte : Espace de Schwartz --- Wikipédia (\url{https://fr.wikipedia.org/wiki/Espace_de_Schwartz})
\item Contexte : Distribution (mathématiques) --- Wikipédia (\url{https://fr.wikipedia.org/wiki/Distribution_(math%C3%A9matiques)}), Théorème de Plancherel --- Wikipédia (\url{https://fr.wikipedia.org/wiki/Th%C3%A9or%C3%A8me_de_Plancherel})
\item Lecture : W. Rudin, \textit{Functional Analysis}, ch. 6--7
\item Lecture : M. Reed et B. Simon, \textit{Methods of Modern Mathematical Physics I: Functional Analysis}, ch. V et IX
\item Lecture : Friedlander \& Joshi, \textit{Introduction to the Theory of Distributions}
\end{refs}

\bigskip

\begin{problem}[{Injections de Sobolev et Rellich--Kondrachov --- Master}]
\textbf{FA-35.} \textbf{Problème.} On dit que $ u \in L^2(0,1) $ admet la \textit{dérivée faible}
$ v \in L^2(0,1) $ si $ \int_0^1 u\varphi' = -\int_0^1 v\varphi $ pour toute
$ \varphi \in C_c^{\infty}(0,1) $, et on note $ H^1(0,1) $ l'espace de ces $ u $, muni de la
norme $ \|u\|_{H^1}^2 = \|u\|_{L^2}^2 + \|u'\|_{L^2}^2 $. Soit $ H_0^1(0,1) $ l'adhérence de
$ C_c^{\infty}(0,1) $ dans $ H^1(0,1) $.
\begin{enumerate}
\item[(a)] Démontrez que $ H^1(0,1) $ est un espace de Hilbert, que toute $ u \in H^1(0,1) $ admet un
représentant continu vérifiant $ u(y) - u(x) = \int_x^y u' $, et déduisez-en les estimations
d'injection
\[ |u(x) - u(y)| \le \|u'\|_{L^2}\,|x-y|^{1/2}, \qquad \|u\|_{\infty} \le C\,\|u\|_{H^1} . \]
\item[(b)] Démontrez l'inégalité de Poincaré sur $ H_0^1(0,1) $ : il existe une constante $ C $ avec
$ \|u\|_{L^2} \le C\,\|u'\|_{L^2} $ pour toute $ u \in H_0^1(0,1) $. Déterminez la meilleure
constante, et identifiez-la comme une valeur propre.
\item[(c)] Démontrez le théorème de Rellich--Kondrachov dans ce cadre : l'injection
$ H_0^1(0,1) \hookrightarrow L^2(0,1) $ est compacte. (Les estimations de (a) rendent une partie
bornée de $ H^1 $ uniformément bornée et équicontinue ; appliquez Arzelà--Ascoli, puis FA-12.)
\item[(d)] Récoltez la conséquence. Étant donné $ f \in L^2(0,1) $, utilisez Lax--Milgram (FA-36) pour
résoudre le problème faible $ \int_0^1 u'v' = \int_0^1 fv $ pour toute $ v \in H_0^1(0,1) $ ;
démontrez que l'opérateur solution $ f \mapsto u $ est compact et autoadjoint sur $ L^2(0,1) $ ; et
concluez de FA-13 que $ L^2(0,1) $ possède une base orthonormée de fonctions propres de
$ -d^2/dx^2 $ avec conditions de Dirichlet au bord, de valeurs propres tendant vers l'infini.
Identifiez-les explicitement et comparez avec FA-05. Énoncez enfin l'inégalité de Sobolev générale
$ W^{1,p}(\mathbb{R}^n) \subset L^{p^*} $ avec $ p^{*} = np/(n-p) $ pour $ 1 \le p < n $, et
expliquez l'argument d'homogénéité $ u_\lambda(x) = u(\lambda x) $ qui force cet exposant et aucun
autre.
\end{enumerate}

\end{problem}

\noindent Sergueï Sobolev a introduit ces espaces dans les années 1930 pour donner un sens aux
solutions de l'équation des ondes qui ne sont pas dérivables, et les théorèmes d'injection sont le taux de
change entre eux et la régularité classique : combien de dérivées faibles dans $ L^p $ achètent combien
de continuité ou d'intégrabilité. Franz Rellich a démontré l'énoncé de compacité pour $ H^1 $ en 1930 et
Vladimir Kondrachov l'a étendu en 1945 ; c'est le théorème de compacité le plus utilisé en équations aux
dérivées partielles, parce que la méthode directe produit des limites faibles (FA-14) et que Rellich est ce
qui les transforme en limites fortes. La question (d) est tout l'appareillage de cette page agissant d'un
seul coup --- Lax--Milgram pour l'existence, Rellich pour la compacité, le théorème spectral de FA-13 pour la
diagonalisation --- et elle se termine là où FA-05 commençait, avec $ \sin(n\pi x) $ et la corde
vibrante.

\begin{refs}
\item Contexte : Espace de Sobolev --- Wikipédia (\url{https://fr.wikipedia.org/wiki/Espace_de_Sobolev})
\item Contexte : Injections de Sobolev --- Wikipédia (\url{https://fr.wikipedia.org/wiki/Injections_de_Sobolev}), Théorème de Rellich--Kondrachov --- Wikipédia (\url{https://fr.wikipedia.org/wiki/Th%C3%A9or%C3%A8me_de_Rellich})
\item Lecture : H. Brezis, \textit{Analyse fonctionnelle : théorie et applications}, ch. 8--9
\item Lecture : L. C. Evans, \textit{Partial Differential Equations}, ch. 5
\end{refs}

\bigskip

\begin{problem}[{Le lemme de Lax--Milgram --- Master, short}]
\textbf{FA-36.} \textbf{Problème.} Soit $ H $ un espace de Hilbert réel et soit
$ a : H \times H \to \mathbb{R} $ bilinéaire, continue ($ |a(u,v)| \le C\|u\|\,\|v\| $) et coercive
($ a(u,u) \ge \alpha \|u\|^2 $ pour un $ \alpha > 0 $), sans hypothèse de symétrie. Démontrez que
pour tout $ f \in H^{*} $ il existe exactement un $ u \in H $ avec $ a(u,v) = f(v) $ pour tout
$ v \in H $, et que $ \|u\| \le \|f\|/\alpha $.
\end{problem}

\noindent Peter Lax et Arthur Milgram ont publié ce résultat en 1954, dans un article sur les
équations paraboliques, et c'est le théorème d'existence de trait pour les problèmes elliptiques
linéaires : la symétrie ayant sauté, le théorème de représentation de Riesz de FA-03(d) survit encore sous
la forme ci-dessus, le théorème du point fixe de Banach (FA-07) ou la série de Neumann (FA-20) fournissant
la surjectivité manquante. Lorsque $ a $ est symétrique, la solution minimise l'énergie
$ \tfrac{1}{2}a(v,v) - f(v) $, si bien que le lemme est le principe de Dirichlet réhabilité --- le
principe même que Weierstrass avait démoli et que Hilbert a réparé. Combiné au lemme de Céa, il est aussi
la licence théorique de la méthode des éléments finis : la solution de Galerkine dans n'importe quel
sous-espace fermé est quasi optimale, avec une constante $ C/\alpha $.

\begin{refs}
\item Contexte : Formulation faible et théorème de Lax--Milgram --- Wikipédia (\url{https://fr.wikipedia.org/wiki/Formulation_faible})
\item Contexte : Lemme de Céa --- Wikipédia (\url{https://fr.wikipedia.org/wiki/Lemme_de_C%C3%A9a})
\item Lecture : H. Brezis, \textit{Analyse fonctionnelle : théorie et applications}, ch. 5
\end{refs}

\bigskip


\section*{Recherche}

\begin{problem}[{Le problème du sous-espace invariant --- Recherche}]
\textbf{FA-17.} \textbf{Problème.} (Ouvert pour l'espace de Hilbert.) Soit $ H $ un espace de Hilbert complexe
séparable de dimension infinie et $ T $ un opérateur borné sur $ H $.
\begin{enumerate}
\item[(a)] $ T $ doit-il posséder un sous-espace fermé invariant autre que $ \{0\} $ et $ H $ --- un
sous-espace fermé $ M $ avec $ T(M) \subset M $ ?
\item[(b)] Comme première entrée sérieuse dans ce cercle d'idées, démontrez le théorème d'Aronszajn--Smith
selon lequel tout opérateur \textit{compact} possède un sous-espace fermé invariant non trivial, ou le
théorème plus fort de Lomonosov (1973) selon lequel tout opérateur commutant avec un opérateur
compact non nul en possède un.
\end{enumerate}

\end{problem}

\noindent La question remonte à von Neumann dans les années 1930. Pour les espaces de Banach la
réponse est \textit{non} : Per Enflo a construit un contre-exemple (annoncé au milieu des années 1970,
publié dans \textit{Acta Mathematica} en 1987), et Charles Read en a donné un sur l'espace concret
$ \ell^1 $ (1985). À l'autre extrême, Argyros et Haydon (2011) ont bâti un espace de Banach sur lequel
tout opérateur borné est scalaire plus compact, de sorte que tout opérateur y possède des sous-espaces
invariants. Pour l'espace de Hilbert le problème reste ouvert. En mai 2023 Enflo a déposé sur arXiv une
prépublication annonçant une solution positive pour les espaces de Hilbert (révisée en 2024) ; en 2026
cette annonce n'a pas été validée par une publication avec comité de lecture et reste en cours
d'évaluation par la communauté, si bien que le statut honnête est : revendiqué, non confirmé, et le
problème doit encore être traité comme ouvert.

\begin{refs}
\item Contexte : Problème du sous-espace invariant --- Wikipédia (en anglais) (\url{https://en.wikipedia.org/wiki/Invariant_subspace_problem})
\item Article : P. Enflo, \og{} On the invariant subspace problem for Banach spaces \fg{} (1987), \textit{Acta Mathematica} 158
\item Article : S. A. Argyros et R. G. Haydon, \og{} A hereditarily indecomposable L$\infty$-space that solves the scalar-plus-compact problem \fg{} (2011), \textit{Acta Mathematica} 206
\item Article : P. Enflo, \og{} On the invariant subspace problem in Hilbert spaces \fg{} (2023, prépublication), arXiv:2305.15442 (\url{https://arxiv.org/abs/2305.15442})
\end{refs}

\bigskip

\begin{problem}[{Le problème de plongement de Connes --- Recherche}]
\textbf{FA-18.} \textbf{Problème.} (Résolu par la négative, 2020.) Un facteur
$ \mathrm{II}_1 $ est une algèbre de von Neumann de dimension infinie, de centre trivial et portant une
trace : l'espace de probabilité non commutatif des analystes ; le facteur $ \mathrm{II}_1 $ hyperfini
$ R $ est le plus petit d'entre eux, limite d'algèbres de matrices. Connes a demandé en 1976 si tout
facteur $ \mathrm{II}_1 $ à prédual séparable se plonge, en préservant la trace, dans une ultrapuissance
$ R^{\omega} $ de $ R $ --- informellement : tout espace de probabilité non commutatif peut-il être
approché par des matrices ? On sait aujourd'hui que la réponse est \textit{non}. Le problème que nous posons :
parcourir la remarquable chaîne d'équivalences --- plongement de Connes $\Leftrightarrow$ conjecture QWEP de Kirchberg $\Leftrightarrow$
problème de Tsirelson sur les corrélations quantiques --- et rédiger un exposé autonome de la façon dont le
théorème de complexité quantique $ \mathrm{MIP}^{*} = \mathrm{RE} $ de 2020 réfute les trois. Notez que
la réfutation est indirecte : extraire de la démonstration un facteur non plongeable explicite et concret
est un problème de recherche actif.
\end{problem}

\noindent Connes a fait sa remarque en passant dans son article de classification de 1976 aux
\textit{Annals}, et pendant quatre décennies elle a organisé les algèbres d'opérateurs comme l'hypothèse de
Riemann organise la théorie des nombres, avec des dizaines d'énoncés démontrés équivalents. La résolution
est venue d'une direction inattendue : en janvier 2020, Ji, Natarajan, Vidick, Wright et Yuen ont démontré
que la classe MIP* des langages décidables par preuves interactives multiprouveurs intriquées égale RE, la
classe des langages récursivement énumérables --- des prouveurs intriqués peuvent donc vérifier le problème
de l'arrêt. Via le problème de Tsirelson cela donne la non-fermeture de l'ensemble des corrélations
quantiques, et donc une réponse négative à Connes. C'est l'un des franchissements de pont les plus
frappants jamais enregistrés entre informatique et mathématiques pures, et son statut est réglé : l'article
a paru après relecture par les pairs et les équivalences qu'il exploite étaient établies dans la
littérature bien avant lui.

\begin{refs}
\item Contexte : Problème de plongement de Connes --- Wikipédia (en anglais) (\url{https://en.wikipedia.org/wiki/Connes_embedding_problem})
\item Article : Z. Ji, A. Natarajan, T. Vidick, J. Wright, H. Yuen, \og{} MIP* = RE \fg{} (2020), arXiv:2001.04383 (\url{https://arxiv.org/abs/2001.04383})
\item Article : A. Connes, \og{} Classification of injective factors \fg{} (1976), \textit{Annals of Mathematics} 104
\end{refs}

\bigskip

\begin{problem}[{Le problème du quotient séparable --- Recherche, short}]
\textbf{FA-28.} \textbf{Problème.} (Ouvert.) Tout espace de Banach $ X $ de dimension infinie
possède-t-il un sous-espace fermé $ M $ tel que le quotient $ X/M $ soit de dimension infinie et
séparable ?
\end{problem}

\noindent La question est attribuée à Banach et Mazur et reste ouverte depuis les années 1930 ---
un scandale de longévité, vu son air innocent. La version \og{} sous-espace \fg{} est facile (tout espace normé de
dimension infinie contient un sous-espace fermé séparable de dimension infinie) ; ce sont les quotients qui
résistent. La réponse est oui pour les espaces réflexifs, pour les espaces $ C(K) $ et $ L^1 $, et --- par
un théorème d'Argyros, Dodos et Kanellopoulos de 2008 --- pour tout espace dual ; en 2026 le problème général
reste ouvert. Le panorama de Mujica rassemble tout un réseau de formulations équivalentes.

\begin{refs}
\item Contexte : Espace de Banach --- Wikipédia (\url{https://fr.wikipedia.org/wiki/Espace_de_Banach})
\item Article : J. Mujica, \og{} Separable quotients of Banach spaces \fg{} (1997), \textit{Revista Matemática de la Universidad Complutense de Madrid} 10
\end{refs}

\bigskip

\begin{problem}[{La conjecture de Crouzeix --- Recherche, short}]
\textbf{FA-29.} \textbf{Problème.} (Ouvert.) L'\textit{image numérique} d'un opérateur borné $ A $ sur
un espace de Hilbert complexe est
$ W(A) = \{\, \langle Ax, x\rangle : \|x\| = 1 \,\} $. Démontrez ou réfutez la conjecture de Crouzeix :
pour un tel $ A $ et tout polynôme $ p $,
\[ \|p(A)\| \le 2 \sup_{z \in W(A)} |p(z)| . \]
\end{problem}

\noindent Michel Crouzeix a posé la conjecture en 2004 et démontré l'inégalité avec la constante
11,08 en 2007 ; la constante 2 serait optimale, l'égalité ayant déjà lieu pour une matrice nilpotente
$ 2 \times 2 $. Crouzeix et Palencia ont ramené la constante à $ 1 + \sqrt{2} $ en 2017, dans une
démonstration admirée pour sa brièveté, et la conjecture a un enjeu pratique : elle certifierait l'image
numérique, calculable, comme estimateur de $ \|p(A)\| $ en algèbre linéaire numérique. Des prépublications
annonçant des démonstrations complètes sont parues à la mi-2026 ; aucune n'a encore été validée par une
relecture par les pairs, si bien que le statut honnête reste ouvert --- peut-être au bord de la
résolution.

\begin{refs}
\item Contexte : Conjecture de Crouzeix --- Wikipédia (en anglais) (\url{https://en.wikipedia.org/wiki/Crouzeix%27s_conjecture})
\item Contexte : Image numérique d'un opérateur --- Wikipédia (en anglais) (\url{https://en.wikipedia.org/wiki/Numerical_range})
\item Article : M. Crouzeix et C. Palencia, \og{} The numerical range is a $ (1+\sqrt 2) $-spectral set \fg{} (2017), \textit{SIAM Journal on Matrix Analysis and Applications} 38
\item Voir aussi : LIN-21 (\url{pure-linear-algebra.html#LIN-21}) énonce la même conjecture pour les matrices finies, où elle est tout aussi ouverte.
\end{refs}

\bigskip

\begin{problem}[{Le problème de Kadison--Singer --- Recherche}]
\textbf{FA-37.} \textbf{Problème.} (Résolu par l'affirmative, 2013.) Un \textit{état} sur une C*-algèbre est une forme
linéaire positive de norme 1, et il est \textit{pur} s'il n'est pas une combinaison convexe propre de deux
autres états. Fixons une base orthonormée $ (e_n) $ de $ \ell^2 $ et soit $ D \subset B(\ell^2) $
l'algèbre des opérateurs diagonaux dans cette base. Kadison et Singer ont demandé en 1959 : tout état pur
sur $ D $ se prolonge-t-il en un \textit{unique} état sur $ B(\ell^2) $ ?
\begin{enumerate}
\item[(a)] Démontrez que $ D $ est une sous-algèbre abélienne autoadjointe maximale de $ B(\ell^2) $,
*-isomorphe isométriquement à $ \ell^\infty $, et que ses états purs sont exactement les caractères de
$ \ell^\infty $.
\item[(b)] Réglez les états purs faciles : démontrez que pour chaque $ n $, l'état
$ \varphi_n(T) = \langle Te_n, e_n \rangle $ sur $ D $ admet un et un seul prolongement en un état
sur $ B(\ell^2) $, à savoir la même formule. (Montrez d'abord que tout prolongement $ \psi $ vérifie
$ \psi(T) = \psi(P_nTP_n) $, où $ P_n $ est le projecteur sur $ \mathbb{C}e_n $.) Expliquez
pourquoi les états purs restants --- ceux provenant d'ultrafiltres libres --- constituent toute la
difficulté.
\item[(c)] Lisez la conjecture de pavage d'Anderson (1979) et la reformulation en termes de discrépance en
dimension finie de Weaver (2004), et rédigez un exposé soigné de la manière dont un énoncé sur le
partage d'ensembles finis de vecteurs peut équivaloir à un énoncé sur les états de
$ B(\ell^2) $.
\item[(d)] Lisez la démonstration de Marcus, Spielman et Srivastava : familles entrelacées de polynômes, méthode
des polynômes caractéristiques mixtes, et majoration de la plus grande racine qui produit une bonne
partition. Prenez ensuite en charge ce que la démonstration ne donne pas : l'argument montre qu'une
partition convenable existe sans en exhiber une, alors cherchez ce que l'on sait de son calcul
efficace.
\end{enumerate}

\end{problem}

\noindent \textbf{Statut : résolu.} Richard Kadison et Isadore Singer ont posé la question en 1959,
dans le sillage de l'insistance de Dirac selon laquelle un système quantique possède un ensemble complet
d'observables qui commutent et dont la mesure détermine l'état ; ils soupçonnaient que la réponse était non.
En cinquante ans le problème s'est révélé équivalent à un éventail saisissant d'énoncés --- la conjecture de
pavage d'Anderson en théorie des opérateurs, la conjecture de discrépance de Weaver, la conjecture de
Feichtinger sur le partage des repères en traitement du signal --- de sorte qu'une solution quelque part les
résolvait tous. Elle est venue en 2013 de trois chercheurs extérieurs aux algèbres d'opérateurs : Adam
Marcus, Daniel Spielman et Nikhil Srivastava, dont la technique des familles entrelacées venait de produire
des graphes de Ramanujan bipartis de tout degré, et qui ont démontré la conjecture de Weaver et donc une
réponse positive. Le travail a paru dans les \textit{Annals of Mathematics} en 2015 et a reçu le prix Pólya
2014 ; c'est le meilleur exemple récent d'un problème tombant sous une idée importée d'un domaine sans
rapport.

\begin{refs}
\item Contexte : Problème de Kadison--Singer --- Wikipédia (en anglais) (\url{https://en.wikipedia.org/wiki/Kadison%E2%80%93Singer_problem})
\item Article : A. Marcus, D. Spielman, N. Srivastava, \og{} Interlacing families II: Mixed characteristic polynomials and the Kadison--Singer problem \fg{} (2015), \textit{Annals of Mathematics} 182, arXiv:1306.3969 (\url{https://arxiv.org/abs/1306.3969})
\item Lecture : J. B. Conway, \textit{A Course in Functional Analysis}, ch. VIII (états et représentations)
\end{refs}

\bigskip

\begin{problem}[{Les facteurs des groupes libres --- Recherche}]
\textbf{FA-38.} \textbf{Problème.} (Ouvert.) Pour un groupe discret $ G $, soit $ \lambda $ la représentation
régulière gauche sur $ \ell^2(G) $ et soit $ L(G) = \lambda(G)'' \subset B(\ell^2(G)) $ l'algèbre de
von Neumann de groupe qu'elle engendre. On note $ F_n $ le groupe libre à $ n $ générateurs.
\begin{enumerate}
\item[(a)] Démontrez que $ \tau(x) = \langle x\delta_e, \delta_e \rangle $ définit un état tracial fidèle sur
$ L(G) $ : $ \tau(xy) = \tau(yx) $, et $ \tau(x^{*}x) = 0 $ entraîne $ x = 0 $.
\item[(b)] Démontrez que $ L(G) $ a un centre trivial --- est un facteur --- si et seulement si $ G $ a des
classes de conjugaison infinies, et vérifiez que $ F_n $ en a pour $ n \ge 2 $. Concluez que
$ L(F_n) $ est un facteur $ \mathrm{II}_1 $, algèbre de von Neumann de dimension infinie munie d'une
trace, comme en FA-18.
\item[(c)] Démontrez, ou reconstituez à partir de l'article de Murray et von Neumann de 1943, que $ L(F_2) $
n'est pas isomorphe au facteur hyperfini $ R $ : montrez que $ R $ a la propriété $\Gamma$ --- il existe des
unitaires approximativement centraux de trace nulle --- et que $ L(F_2) $ ne l'a pas.
\item[(d)] Voici maintenant la question ouverte : a-t-on $ L(F_2) \cong L(F_3) $ ? Lisez les probabilités
libres de Voiculescu, les facteurs de groupes libres interpolés $ L(F_t) $ de Dykema et Rădulescu, et
leur dichotomie, puis rédigez un exposé expliquant pourquoi les invariants standard échouent --- y compris
le théorème de Voiculescu de 1996 selon lequel les facteurs de groupes libres n'ont pas de sous-algèbre
de Cartan, ce qui a tué la stratégie la plus naturelle pour les séparer.
\end{enumerate}

\end{problem}

\noindent \textbf{Statut : ouvert.} Murray et von Neumann ont construit les premiers exemples de
facteurs $ \mathrm{II}_1 $ dans les années 1930 et se sont aussitôt demandé combien il y en avait ; ils
ont séparé $ L(F_2) $ du facteur hyperfini $ R $ par la propriété $\Gamma$ en 1943, et la question des groupes
libres en est restée là. Il n'est pas seulement inconnu que $ L(F_2) \cong L(F_3) $ --- on ne sait même pas
si deux facteurs de groupes libres quelconques sont isomorphes. Voiculescu a inventé les probabilités libres
dans les années 1980 pour attaquer exactement cela, et elles ont produit une étrange réponse partielle :
Dykema et Rădulescu ont montré indépendamment au début des années 1990 que les facteurs interpolés
$ L(F_t) $, définis pour tout réel $ t > 1 $, obéissent à une dichotomie du tout ou rien --- ou bien
tous les $ L(F_n) $ avec $ 2 \le n \le \infty $ sont isomorphes, ou bien aucun couple ne l'est. La
dimension d'entropie libre a été conçue pour trancher et n'y est pas encore parvenue.

\begin{refs}
\item Contexte : Algèbre de von Neumann --- Wikipédia (\url{https://fr.wikipedia.org/wiki/Alg%C3%A8bre_de_von_Neumann})
\item Contexte : Probabilité libre --- Wikipédia (\url{https://fr.wikipedia.org/wiki/Probabilit%C3%A9_libre})
\item Article : F. J. Murray et J. von Neumann, \og{} On rings of operators. IV \fg{} (1943), \textit{Annals of Mathematics} 44
\item Lecture : Voiculescu, Dykema \& Nica, \textit{Free Random Variables} (CRM Monograph Series, 1992)
\end{refs}

\bigskip

\begin{problem}[{La constante de Grothendieck --- Recherche, short}]
\textbf{FA-39.} \textbf{Problème.} (Ouvert.) L'inégalité de Grothendieck affirme l'existence d'une
constante universelle $ K $ telle que, pour toute matrice réelle $ (m_{ij}) $ vérifiant
$ \left| \sum_{i,j} m_{ij} s_i t_j \right| \le 1 $ pour tous réels
$ s_i, t_j \in [-1, 1] $, et pour tous vecteurs unitaires $ x_i, y_j $ d'un espace de Hilbert réel
quelconque,
\[ \left| \sum_{i,j} m_{ij} \langle x_i, y_j \rangle \right| \;\le\; K . \]
Déterminez la valeur exacte de la plus petite constante possible, la constante de Grothendieck réelle
$ K_G^{\mathbb{R}} $.
\end{problem}

\noindent \textbf{Statut : ouvert.} Grothendieck a démontré l'inégalité dans son
\textit{Résumé de la théorie métrique des produits tensoriels topologiques} de 1953, article si en avance sur
son public qu'il fut ignoré quinze ans, jusqu'à ce que Lindenstrauss et Pełczyński le réécrivent dans le
langage de la théorie des espaces de Banach. On sait seulement que la constante se situe entre environ
1,6769, minoration de Davie et de Reeds, et la majoration de Krivine de 1979,
$ \pi / (2\log(1+\sqrt{2})) \approx 1{,}7822 $, dont Krivine conjecturait qu'elle était exacte ;
Braverman, Makarychev, Makarychev et Naor l'ont réfuté en 2011 en montrant que la vraie valeur est
strictement plus petite, sans produire de nouvelle majoration explicite. L'inégalité a depuis complètement
échappé à l'analyse fonctionnelle : elle contrôle la qualité des relaxations semi-définies de MAX-CUT en
informatique théorique, et elle borne les corrélations réalisables par des systèmes quantiques
intriqués.

\begin{refs}
\item Contexte : Inégalité de Grothendieck --- Wikipédia (en anglais) (\url{https://en.wikipedia.org/wiki/Grothendieck_inequality})
\item Article : A. Grothendieck, \og{} Résumé de la théorie métrique des produits tensoriels topologiques \fg{} (1953), \textit{Boletim da Sociedade de Matemática de São Paulo} 8
\item Lecture : Pisier, \textit{Similarity Problems and Completely Bounded Maps}
\end{refs}

\bigskip


\section*{Pour aller plus loin}


\vfill
\noindent\emph{Hors de la Caverne} --- des probl\`emes, pas de r\'eponses.
\end{document}
